% concatenated from cvar_bounds_arxiv.tex
%File: anonymous-submission-latex-2026.tex
\documentclass[letterpaper]{article} % DO NOT CHANGE THIS
\usepackage{aaai2026}  % DO NOT CHANGE THIS
\usepackage{times}  % DO NOT CHANGE THIS
\usepackage{helvet}  % DO NOT CHANGE THIS
\usepackage{courier}  % DO NOT CHANGE THIS
\usepackage[hyphens]{url}  % DO NOT CHANGE THIS
\usepackage{graphicx} % DO NOT CHANGE THIS
\urlstyle{rm} % DO NOT CHANGE THIS
\def\UrlFont{\rm}  % DO NOT CHANGE THIS
\usepackage{natbib}  % DO NOT CHANGE THIS AND DO NOT ADD ANY OPTIONS TO IT
\usepackage{caption} % DO NOT CHANGE THIS AND DO NOT ADD ANY OPTIONS TO IT
\frenchspacing  % DO NOT CHANGE THIS
\setlength{\pdfpagewidth}{8.5in} % DO NOT CHANGE THIS
\setlength{\pdfpageheight}{11in} % DO NOT CHANGE THIS
%
% These are recommended to typeset algorithms but not required. See the subsubsection on algorithms. Remove them if you don't have algorithms in your paper.
\usepackage{algorithm}
\usepackage{algorithmic}

%
% These are are recommended to typeset listings but not required. See the subsubsection on listing. Remove this block if you don't have listings in your paper.
\usepackage{newfloat}
\usepackage{listings}
\DeclareCaptionStyle{ruled}{labelfont=normalfont,labelsep=colon,strut=off} % DO NOT CHANGE THIS
\lstset{%
	basicstyle={\footnotesize\ttfamily},% footnotesize acceptable for monospace
	numbers=left,numberstyle=\footnotesize,xleftmargin=2em,% show line numbers, remove this entire line if you don't want the numbers.
	aboveskip=0pt,belowskip=0pt,%
	showstringspaces=false,tabsize=2,breaklines=true}
\floatstyle{ruled}
\newfloat{listing}{tb}{lst}{}
\floatname{listing}{Listing}
%
% Keep the \pdfinfo as shown here. There's no need
% for you to add the /Title and /Author tags.
\pdfinfo{
	/TemplateVersion (2026.1)
}

\newtheorem{theorem}{Theorem}[section]
\newtheorem{lemma}{Lemma}[section]
\newtheorem{definition}{Definition}
\newtheorem{corollary}[theorem]{Corollary}

\newenvironment{proof}{\noindent\textit{Proof.}}{\hfill$\square$}
\newcommand{\theoremlabel}{\labelcref} % Define a label for your theorem

\newcommand{\KP}[1]{\textcolor{red}{#1}}
\newcommand{\VI}[1]{\textcolor{blue}{#1}}

% My packages
\usepackage{amsmath}   % for \overset, \text
\usepackage{amssymb}   % for \triangleq, \mathbb
\usepackage{amsfonts}  % for \mathbb (optional; amssymb usually includes it)
\usepackage{xcolor}
\usepackage{subfig}


% DISALLOWED PACKAGES
% \usepackage{authblk} -- This package is specifically forbidden
% \usepackage{balance} -- This package is specifically forbidden
% \usepackage{color (if used in text)
	% \usepackage{CJK} -- This package is specifically forbidden
	% \usepackage{float} -- This package is specifically forbidden
	% \usepackage{flushend} -- This package is specifically forbidden
	% \usepackage{fontenc} -- This package is specifically forbidden
	% \usepackage{fullpage} -- This package is specifically forbidden
	% \usepackage{geometry} -- This package is specifically forbidden
	% \usepackage{grffile} -- This package is specifically forbidden
	% \usepackage{hyperref} -- This package is specifically forbidden
	% \usepackage{navigator} -- This package is specifically forbidden
	% (or any other package that embeds links such as navigator or hyperref)
	% \indentfirst} -- This package is specifically forbidden
% \layout} -- This package is specifically forbidden
% \multicol} -- This package is specifically forbidden
% \nameref} -- This package is specifically forbidden
% \usepackage{savetrees} -- This package is specifically forbidden
% \usepackage{setspace} -- This package is specifically forbidden
% \usepackage{stfloats} -- This package is specifically forbidden
% \usepackage{tabu} -- This package is specifically forbidden
% \usepackage{titlesec} -- This package is specifically forbidden
% \usepackage{tocbibind} -- This package is specifically forbidden
% \usepackage{ulem} -- This package is specifically forbidden
% \usepackage{wrapfig} -- This package is specifically forbidden
% DISALLOWED COMMANDS
% \nocopyright -- Your paper will not be published if you use this command
% \addtolength -- This command may not be used
% \balance -- This command may not be used
% \baselinestretch -- Your paper will not be published if you use this command
% \clearpage -- No page breaks of any kind may be used for the final version of your paper
% \columnsep -- This command may not be used
% \newpage -- No page breaks of any kind may be used for the final version of your paper
% \pagebreak -- No page breaks of any kind may be used for the final version of your paperr
% \pagestyle -- This command may not be used
% \tiny -- This is not an acceptable font size.
% \vspace{- -- No negative value may be used in proximity of a caption, figure, table, section, subsection, subsubsection, or reference
% \vskip{- -- No negative value may be used to alter spacing above or below a caption, figure, table, section, subsection, subsubsection, or reference

\setcounter{secnumdepth}{0} %May be changed to 1 or 2 if section numbers are desired.

% The file aaai2026.sty is the style file for AAAI Press
% proceedings, working notes, and technical reports.
%

% Title

% Your title must be in mixed case, not sentence case.
% That means all verbs (including short verbs like be, is, using,and go),
% nouns, adverbs, adjectives should be capitalized, including both words in hyphenated terms, while
% articles, conjunctions, and prepositions are lower case unless they
% directly follow a colon or long dash
%\title{Bounding Conditional Value-at-Risk via Distributional Discrepancies}
%\title{Bounding Conditional Value-at-Risk via Cumulative Distributions Discrepancies}
\title{Bounding Conditional Value-at-Risk via Auxiliary Distributions with Bounded Discrepancies}
\author{
%Authors
% All authors must be in the same font size and format.
Yaacov Pariente$^1$, Vadim Indelman$^{2,3}$\\
}
\affiliations{
	%Afiliations
	\textsuperscript{\rm 1}Faculty of Mathematics\\%, Technion - Israel Institute of Technology\\
	\textsuperscript{\rm 2}Stephen B. Klein Faculty of Aerospace Engineering\\%, Technion - Israel Institute of Technology\\
	\textsuperscript{\rm 3}Faculty of Data and Decision Sciences\\%, Technion - Israel Institute of Technology\\
	% If you have multiple authors and multiple affiliations
	% use superscripts in text and roman font to identify them.
	% For example,
	
	% Sunil Issar\textsuperscript{\rm 2}, 
	% J. Scott Penberthy\textsuperscript{\rm 3}, 
	% George Ferguson\textsuperscript{\rm 4},
	% Hans Guesgen\textsuperscript{\rm 5}
	% Note that the comma should be placed after the superscript
	% email address must be in roman text type, not monospace or sans serif
	Technion - Israel Institute of Technology\\
	yaacovp@campus.technion.ac.il, vadim.indelman@technion.ac.il
	%
	% See more examples next
}



\iffalse
%Example, Multiple Authors, ->> remove \iffalse,\fi and place them surrounding AAAI title to use it
\title{My Publication Title --- Multiple Authors}
\author {
% Authors
First Author Name\textsuperscript{\rm 1},
Second Author Name\textsuperscript{\rm 2},
Third Author Name\textsuperscript{\rm 1}
}
\affiliations {
% Affiliations
\textsuperscript{\rm 1}Affiliation 1\\
\textsuperscript{\rm 2}Affiliation 2\\
firstAuthor@affiliation1.com, secondAuthor@affilation2.com, thirdAuthor@affiliation1.com
}
\fi


% REMOVE THIS: bibentry
% This is only needed to show inline citations in the guidelines document. You should not need it and can safely delete it.
\usepackage{bibentry}
% END REMOVE bibentry

\begin{document}

\maketitle

\begin{abstract}
	In this paper, we develop a theoretical framework for bounding the CVaR of a random variable $X$ using another related random variable $Y$, under assumptions on their cumulative and density functions. Our results yield practical tools for approximating $\operatorname{CVaR}_\alpha(X)$ when direct information about $X$ is limited or sampling is computationally expensive, by exploiting a more tractable or observable random variable $Y$. Moreover, the derived bounds provide interpretable concentration inequalities that quantify how the tail risk of $X$ can be controlled via $Y$.
\end{abstract}


\section{Introduction}
Conditional Value-at-Risk (CVaR) has emerged as a critical risk measure in modern artificial intelligence, particularly in settings where decision-making under uncertainty and rare but severe events must be rigorously addressed. Unlike traditional expected value objectives, which can overlook tail risks, CVaR explicitly quantifies the expected loss in the worst-case fraction of scenarios, providing a principled framework for robust and risk-sensitive learning and planning. This property makes CVaR especially valuable in reinforcement learning, safe control, and robust optimization, where agents must balance performance with safety guarantees in the presence of model misspecification, adversarial disturbances, or highly uncertain environments. As AI systems increasingly operate in high-stakes domains—from autonomous driving to financial trading—the ability to reason about and control tail risks via CVaR is becoming an essential component of reliable and trustworthy AI.

The Conditional Value at Risk (CVaR) \cite{rockafellar2000optimization} is a widely used risk measure that facilitates the optimization of the upper tail of the cost's distribution. The dual representation of CVaR \cite{artzner1999coherent} enables its interpretation as the worst-case expectation of the cost \cite{chow2015risk}, thereby motivating its use for risk-averse decision making. CVaR is also a coherent risk measure with desirable properties for safe planning \cite{majumdar2020should}, and its estimators have performance guarantees that ensure their reliability in practice \cite{brown2007large,pmlr-v97-thomas19a}.

In this paper, we derive bounds on the CVaR of a random variable $X$, given another random variable $Y$, where the relationship between $X$ and $Y$ is characterized in terms of their cumulative distribution functions (CDFs) and probability density functions (PDFs). These bounds have three key contributions in practice:
\begin{enumerate}
	\item They enable the approximation of $\operatorname{CVaR}_\alpha(X)$ when direct access to $X$ is unavailable, by leveraging a related distribution $Y$ for which information is accessible.
	\item They provide a means to approximate $\operatorname{CVaR}_\alpha(X)$ when sampling from $X$ is computationally costly, by instead using a distribution $Y$ that is less expensive to simulate or estimate.
	\item They lead to interpretable concentration inequalities, encompassing existing concentration bounds for $\operatorname{CVaR}_\alpha(X)$ as a special case. Specifically, they characterize the CVaR of $X$ through the CVaR of $Y$ with an appropriately adjusted confidence level.
\end{enumerate}
Proofs that do not appear in the main text are provided in the supplemental material.

\section{Preliminaries}
\addtocounter{section}{1}
Let $X$ be a random variable defined on a probability space $(\Omega, \mathcal{F}, P)$, where $\mathcal{F} = 2^\Omega$ is the $\sigma$-algebra, and $P : \mathcal{F} \to [0,1]$ is a probability measure. Assume further that $E|X| < \infty$. We denote the cumulative density function (CDF) of the random variable $X$ by $F_X(x)=P(X\leq x)$. The value at risk at confidence level $\alpha \in (0,1)$ is the $1-\alpha$ quantile of $X$, i.e \begin{equation}
	VaR_\alpha (X) \triangleq \inf\{x\in \mathbb{R}:F(x) > 1-\alpha\}
\end{equation}
For simplicity we denote $q_\alpha^X=VaR_\alpha(X)$ or $q_\alpha^F=VaR_\alpha(X)$. The conditional value at risk (CVaR) at confidence level $\alpha$ is defined as \cite{rockafellar2000optimization}
\begin{equation}
	CVaR_\alpha(X) \triangleq \inf_{w\in \mathbb{R}}\{w+\frac{1}{\alpha}\mathbb{E}[(X-w)^+]|w\in \mathbb{R}\},
\end{equation}
where $(x)^+=\max{(x,0)}$. For a smooth $F$, it holds that \cite{pflug2000some}
\begin{equation}\label{def:cvar_as_conditional_expectation}
	CVaR_\alpha(X)=\mathbb{E}[X|X>VaR_\alpha(X)]=\frac{1}{\alpha}\int_{1-\alpha}^1 F^{-1}(v)dv.
\end{equation}
Let $X_i \overset{\text{iid}}{\sim} F$ for $i\in\{1,\dots,n\}$. Denote by
\begin{equation}\label{eq:brown_cvar_estimator}
	\hat{C}_\alpha(X) \triangleq \hat{C}_\alpha(\{X_i\}^{n}_{i=1}) \triangleq \inf_{x\in \mathbb{R}} \Bigl{\{}x+\frac{1}{n\alpha} \sum_{i=1}^n(X_i-x)^+\Bigr{\}}
\end{equation}
the estimate of $CVaR_\alpha(X)$ \cite{brown2007large}. Theorem \ref{thm:brown_bounds}, 
that bounds the deviation of the estimated CVaR and the true CVaR, was proved in \cite{brown2007large}.
\begin{theorem}\label{thm:brown_bounds}
	If $\text{supp}(X) \subseteq [a, b]$ and $X$ has a continuous distribution function, then for any $\delta \in (0, 1]$, 
	\begin{equation}
		P\Bigl{(}CVaR_\alpha(X)-\hat{C}_\alpha(X)>(b-a)\sqrt{\frac{5ln(3/\delta)}{\alpha n}}\Bigr{)}\leq \delta,
	\end{equation}
	\begin{equation}
		P\Bigl{(}CVaR_\alpha(X)-\hat{C}_\alpha(X)<-\frac{(b-a)}{\alpha}\sqrt{\frac{ln(1/\delta)}{2n}}\Bigr{)}\leq \delta.
	\end{equation}
\end{theorem}

\eqref{eq:brown_cvar_estimator} can be expressed as
\begin{equation}\label{eq:cvar_estimator_sorted_sample}
	\hat{C}_\alpha(X)=X^{(n)}-\frac{1}{\alpha}\sum_{i=1}^{n-1} (X^{(i)}-X^{(i-1)})\Bigl{(} \frac{i}{n}-(1-\alpha) \Bigr{)}^+,
\end{equation}
where $X^{(i)}$ is the $i$th order statistic of $X_1,\dots,X_n$ in ascending order \cite{pmlr-v97-thomas19a}. The results presented in Theorems~\ref{thm:thomas_upper_bound} and~\ref{thm:thomas_lower_bound}, following the work of \cite{pmlr-v97-thomas19a}, yield tighter bounds on the CVaR compared to those established by \cite{brown2007large}.
\begin{theorem}\label{thm:thomas_upper_bound}
	If $X_1, \ldots, X_n$ are independent and identically distributed random variables and $\Pr(X_1 \le b) = 1$ for some finite $b$, then for any $\delta \in (0, 0.5]$,
	\begin{equation}
		\begin{aligned}
			&\Pr\Bigg(
			\operatorname{CVaR}_\alpha(X_1)
			\le Z_{n+1}
			- \frac{1}{\alpha} \sum_{i=1}^n (Z_{i+1} - Z_i) \\
			&\times 
			\Bigg(
			\frac{i}{n} - \sqrt{\frac{\ln(1/\delta)}{2n}}
			- (1 - \alpha)\Bigg)^{+}
			\Bigg) \ge 1 - \delta,
		\end{aligned}
	\end{equation}
	where $Z_1, \ldots, Z_n$ are the order statistics (i.e., $X_1, \ldots, X_n$ sorted in ascending order), $Z_{n+1} = b$, and $x^+ \triangleq  \max\{0, x\}$ for all $x \in \mathbb{R}$.
\end{theorem}

\begin{theorem}\label{thm:thomas_lower_bound}
	If $X_1, \ldots, X_n$ are independent and identically distributed random variables and $\Pr(X_1 \ge a) = 1$ for some finite $a$, then for any $\delta \in (0, 0.5]$,
	\begin{equation}
		\begin{aligned}
			&\Pr\Bigg(
			\operatorname{CVaR}_\alpha(X_1)
			\ge Z_n
			- \frac{1}{\alpha} \sum_{i=0}^{n-1} (Z_{i+1} - Z_i) \\
			&\times \Big(
			\min \Big\{
			1,\, \frac{i}{n} + \sqrt{\frac{\ln(1/\delta)}{2n}}
			\Big\}
			- (1 - \alpha) \Big)^{+}
			\Bigg) \ge 1 - \delta,
		\end{aligned}
	\end{equation}
	where $Z_1, \ldots, Z_n$ are the order statistics (i.e., $X_1, \ldots, X_n$ sorted in ascending order), $Z_0 = a$, and $x^+ \triangleq \max\{0, x\}$ for all $x \in \mathbb{R}$.
\end{theorem}

\section{Problem Formulation}
Let $X$ and $Y$ be two random variables. This paper seeks to establish upper and lower bounding functions, $f^U$ and $f^L$, which take as input the random variable $Y$ and a confidence level $\alpha$, and satisfy
\begin{equation}
	f^L(Y, \alpha)\leq CVaR_\alpha(X) \leq f^U(Y,\alpha).
\end{equation} 

\section{Theoretical CVaR Bounds}
\addtocounter{section}{1}
In this section, we establish bounds for the CVaR of the random variable $X$ by leveraging an auxiliary random variable $Y$. Two forms of distributional relationships between their respective CDFs are considered: a uniform bound and a non-uniform bound, as illustrated in Figure~\ref{fig:eps_bound} and Figure~\ref{fig:g_bound}. These results provide the theoretical basis for subsequent sections, in which the derived bounds are applied to estimate the CVaR of random variables that are either computationally intractable or prohibitively expensive to sample directly.

Theorem~\ref{thm:cvar_bound_v2} bounds $\mathrm{CVaR}_\alpha(X)$ in terms of $\mathrm{CVaR}_\alpha(Y)$ under the sole condition that the cumulative distribution functions of $X$ and $Y$ differ by at most a known uniform bound. This representation enhances the interpretability of the bound and constitutes a novel aspect of the result, made possible by framing the bounding problem in terms of distributional discrepancies. In the next section, we demonstrate that the bounds established in \cite{pmlr-v97-thomas19a} (theorems \ref{thm:thomas_upper_bound} and \ref{thm:thomas_lower_bound}) arise as a special case of Theorem \ref{thm:cvar_bound_v2}, thereby providing an interpretation for existing CVaR bounds that are otherwise difficult to interpret.
\begin{figure}
	\centering
	\includegraphics[width=0.8\linewidth]{F_Y_bounds_plot.png}
	\caption{Illustration of the bounds on $F_X$.}
	\label{fig:CDF_bounds}
\end{figure}
%\begin{theorem}\label{thm:cvar_bound_v2}
%	Let $X\in[a_X, b_X]$ and $Y\in [a_Y,b_Y]$ be random variables for $a_X, a_Y, b_X, b_Y \in \mathbb{R}$, and denote $a=max(a_X,a_Y),b=min(b_X,b_Y)$. Let $\epsilon \in [0, 1]$. 
%	\begin{enumerate}
%		\item If $\forall z\in \mathbb{R},F_Y(z)-F_X(z)\leq \epsilon$ and $\alpha > \epsilon$ then 
%		\begin{equation}\label{eq:general_cvar_upper_bound_1}
%			CVaR_\alpha(X) \leq \frac{\epsilon}{\alpha}max(b_X,b_Y) + (1 - \frac{\epsilon}{\alpha})CVaR_{\alpha-\epsilon}(Y).
%		\end{equation}
%		
%		\item If $\forall z\in \mathbb{R},F_Y(z)-F_X(z)\leq \epsilon$ and $\alpha \leq \epsilon$ then $CVaR_\alpha(X)\leq max(b_X,b_Y)$.
%		
%		\item If $\forall z\in \mathbb{R},F_X(z)-F_Y(z)\leq \epsilon$ and $\alpha+\epsilon \leq 1$ then 
%		\begin{equation}\label{eq:general_cvar_lower_bound_1}
%			CVaR_{\alpha}(X)\geq (1+\frac{\epsilon}{\alpha})CVaR_{\alpha+\epsilon}(Y)-\frac{\epsilon}{\alpha} CVaR_{\epsilon}(Y)
%		\end{equation}
%		
%		\item If $\forall z\in \mathbb{R},F_X(z)-F_Y(z)\leq \epsilon$ and $\alpha+\epsilon > 1$ then 
%		\begin{equation}
%			CVaR_\alpha(X) \geq \mathbb{E}[Y] - \epsilon CVaR_\epsilon(Y) + (\alpha+\epsilon-1)a_{min}
%		\end{equation}
%	\end{enumerate}
%\end{theorem}
\begin{theorem}\label{thm:cvar_bound_v2}
	Let $X$ and $Y$ be random variables and $\epsilon \in [0, 1]$. 
	\begin{enumerate}
		\item \textbf{Upper bound:} assume that $P(X\leq b_X)=1, P(Y\leq b_Y)=1$ and $\forall z\in \mathbb{R},F_Y(z)-F_X(z)\leq \epsilon$, then 
		\begin{enumerate}
			\item If $\alpha > \epsilon$, \begin{equation}\label{eq:general_cvar_upper_bound_1}
				CVaR_\alpha(X) \leq \frac{\epsilon}{\alpha}max(b_X,b_Y) + (1 - \frac{\epsilon}{\alpha})CVaR_{\alpha-\epsilon}(Y).
			\end{equation}
			
			\item If $\alpha \leq \epsilon$, then $CVaR_\alpha(X)\leq max(b_X,b_Y)$.
		\end{enumerate}
		
		\item \textbf{Lower Bound:} If $P(X\geq a_X)=1, P(Y\geq a_Y)=1$ and $\forall z\in \mathbb{R},F_X(z)-F_Y(z)\leq \epsilon$, then
		\begin{enumerate}
			\item If $\alpha+\epsilon \leq 1$, \begin{equation}\label{eq:general_cvar_lower_bound_1}
				CVaR_{\alpha}(X)\geq (1+\frac{\epsilon}{\alpha})CVaR_{\alpha+\epsilon}(Y)-\frac{\epsilon}{\alpha} CVaR_{\epsilon}(Y)
			\end{equation}
			
			\item If $\alpha+\epsilon > 1$ and $a_{min} = min(a_X, a_Y)$, \begin{equation}
				CVaR_\alpha(X) \geq \mathbb{E}[Y] - \epsilon CVaR_\epsilon(Y) + (\alpha+\epsilon-1)a_{min}
			\end{equation}
		\end{enumerate}
	\end{enumerate}
\end{theorem}
\begin{proof}
	Due to space constraints, we provide only a proof sketch here; the complete proof can be found in the supplementary material. 
	
	The strategy of the proof is to construct two distributions derived from $F_Y$, denoted $F_Y^L$ and $F_Y^U$, such that $F_X$ is first-order stochastically dominated by $F_Y^U$ and first-order stochastically dominates $F_Y^L$; that is, $\forall x\in \mathbb{R},F_Y^L(x) \geq F_X \geq F_Y^U(x)$. Consequently, since CVaR is a coherent risk measure, it follows that
	\begin{equation}
		CVaR_\alpha^{F_Y^L} \leq CVaR_\alpha^{F_X} \leq CVaR_\alpha^{F_Y^U},
	\end{equation}
	where $CVaR_\alpha^{F_Y^L}$, $CVaR_\alpha^{F_X}$, and $CVaR_\alpha^{F_Y^U}$ denote the CVaR at level $\alpha$ corresponding to the distributions $F_Y^L$, $F_X$, and $F_Y^U$, respectively. 
	
	Let $a_{\min} = \min(a_X, a_Y)$ and $b_{\min} = \min(b_X, b_Y)$, and assume $P(a_X\leq X \leq b_X)=1, P(a_Y\leq Y \leq b_Y)=1$. Define the interval $[b_{\min}, b_X)$ to be $\emptyset$ if $b_{\min} = b_X$, and equal to $[b_Y, b_X)$ otherwise. Analogously, define $[a_{\min}, a_X)$ in the same manner. We then define upper and lower bounds for $F_Y$ as follows (see Figure \ref{fig:CDF_bounds}): 
	\begin{equation}
		F_Y^U(y) = 
		\begin{cases}
			0 & y < \max(a_X, q_{1-\epsilon}^Y) \\
			\min(F_Y(y) - \epsilon, 1 - \epsilon) 
			& y \in [\max(a_X, q_{1-\epsilon}^Y),\\
			& , b_{\min}) \\
			1 - \epsilon & y \in [b_{\min}, b_X) \\
			1 & y \geq \max(b_X, b_Y)
		\end{cases}
	\end{equation}
	\begin{equation}
		F_Y^L(y)=\begin{cases}
			0 & y<a_{min}  \\
			\epsilon & y\in [a_{min}, a_Y) \\
			min(F_Y(y)+\epsilon, 1) & y\in [a_Y, min(q_{\epsilon}^Y, b_X)) \\
			1 & y\geq min(q_\epsilon^Y, b_X).
		\end{cases}
	\end{equation}
	As a first step, it is necessary to verify that $F_Y^L$ and $F_Y^U$ are valid cumulative distribution functions and that they satisfy $F_Y^L \leq F_X \leq F_Y^U$. This verification is deferred to the supplementary material. Figure \ref{fig:CDF_bounds} illustrates the corresponding upper and lower bounding CDFs.
	
	\textbf{Upper bound for $CVaR_\alpha^{F_Y^U}$:} From \cite{acerbi2002coherence}, CVaR is equal to an integral of the VaR
	\begin{equation}
		\begin{aligned}
			&CVaR_\alpha^{F_Y^U}=\frac{1}{\alpha}\int_{1-\alpha}^1 \inf \{y\in \mathbb{R}:F_Y^U(y)\geq \tau\} d\tau \\
			&=\frac{1}{\alpha}\int_{1-\alpha+\epsilon}^{1+\epsilon} \inf \{y\in \mathbb{R}:F_Y^U(y)\geq \tau -\epsilon\} d\tau
		\end{aligned}
	\end{equation}  
	If $\alpha \leq \epsilon$, then $1 - \alpha + \epsilon \geq 1$, rendering the bound in the preceding equation trivial, as it attains the maximum value of the support of both $X$ and $Y$.
	\begin{equation}\label{eq:helper1}
		\begin{aligned}
			&\frac{1}{\alpha}\int_{1-\alpha+\epsilon}^{1+\epsilon} \inf \{y\in \mathbb{R}:F_Y^U(y)\geq \tau -\epsilon\} \\
			&\leq \frac{1}{\alpha}\int_{1-\alpha+\epsilon}^{1+\epsilon} max(b_X, b_Y)d\tau = max(b_X, b_Y).
		\end{aligned}
	\end{equation}
	That is, $CVaR_\alpha(X)\leq max(b_X, b_Y)$.
	
	If $\alpha > \epsilon$, then $1 - \alpha + \epsilon < 0$, and the integral may be decomposed into one term that is trivially bounded by the maximum of the support and another term that can be computed explicitly
	\begin{equation}
		\begin{aligned}
			&CVaR_\alpha^{F_Y^U}=\frac{1}{\alpha}\int_{1-\alpha+\epsilon}^{1+\epsilon} \inf \{y\in \mathbb{R}:F_Y^U(y)\geq \tau -\epsilon\} d\tau \\
			&=\frac{1}{\alpha}[\underbrace{\int_{1}^{1+\epsilon} \inf \{y\in \mathbb{R}:F_Y^U(y)\geq \tau -\epsilon\} d\tau}_{\triangleq A_1} \\
			&+\underbrace{\int_{1-\alpha+\epsilon}^{1} \inf \{y\in \mathbb{R}:F_Y^U(y)\geq \tau -\epsilon\} d\tau}_{\triangleq A_2}]
		\end{aligned}
	\end{equation}
	The term $A_1$ is bounded, in a manner analogous to \eqref{eq:helper1}, by $\epsilon \max(b_X, b_Y)$, which is essentially the tightest bound attainable given the definition of $F_Y^U$. The term $A_2$ can be expressed in terms of the CVaR of $Y$, evaluated at a shifted confidence level. Specifically, the variable $\tau - \epsilon$ in $A_2$ lies within the interval $[1 - \alpha, 1 - \epsilon]$. Over this range, for all $y \in [\max(a_X, q_{1-\epsilon}^Y), b_{\min}]$, the inequality $F_Y^U(y) \leq F_Y(y) - \epsilon$ holds. This follows because if $q_{1-\epsilon}^Y \geq a_X$, then by definition $F_Y^U(y) = min(F_Y(y) - \epsilon, 1-\epsilon)$, and if $a_X > q_{1-\epsilon}^Y$, then $F_Y^U(y)=0 \leq F_Y(y) - \epsilon$ for $y\in [q_{1-\epsilon}^Y, a_X]$ and again $F_Y^U(y) = min(F_Y(y) - \epsilon, 1-\epsilon)$ for $y\in [a_X, b_{min}]$.
	\begin{equation}
		\begin{aligned}
			&A_2 \leq \int_{1-\alpha+\epsilon}^1 \inf \{y\in \mathbb{R}:F_Y(y) - \epsilon \geq \tau -\epsilon\}d\tau \\
			&=\int_{1-(\alpha-\epsilon)}^1 \inf \{y\in \mathbb{R}:F_Y(y) \geq \tau\}d\tau \\
			&=(\alpha-\epsilon)CVaR_{\alpha-\epsilon}(Y).
		\end{aligned}
	\end{equation}
	Note that the confidence level is valid, as $\alpha, \epsilon \in (0,1)$ and $\alpha > \epsilon$ imply $\alpha-\epsilon \in (0, 1)$. By combining the previous two expressions, we obtain a single bound:
	\begin{equation}
		CVaR_\alpha^{F_Y^U} \leq \epsilon\frac{max(b_X,b_Y)}{\alpha} + (1 - \frac{\epsilon}{\alpha})CVaR_{\epsilon-\alpha}(Y).
	\end{equation}
	In the case where $a_X = -\infty$, we have $\max(a_X, q_{1-\epsilon}^Y) = q_{1-\epsilon}^Y$, and the definition of $F_Y^U$ no longer involves $a_X$. Therefore, the inequality remains valid in the case of $a_X = -\infty$.
		
	\textbf{Lower bound for $CVaR_\alpha^{F_Y^L}$:} The detailed proof is provided in the supplementary material. The proof follows an approach analogous to that of the upper bound: specifically, the integral defining $CVaR_\alpha^{F_Y^L}$ is decomposed into terms that exhibit the structure of a CVaR expression.
\end{proof}

The parameter $\epsilon$ quantifies the discrepancy between the random variables $X$ and $Y$, and is defined as a uniform bound on the difference between their cumulative distribution functions. By construction, $\epsilon$ takes values in the interval $[0,1]$. In the limiting case where $\epsilon = 1$, $Y$ provides no information about $X$. In this setting, the bounds provided by Theorem~\ref{thm:cvar_bound_v2} reduce to trivial bounds on $\mathrm{CVaR}_\alpha(X)$ involving the essential supremum of both $X$ and $Y$, rendering them ineffective for practical use. When $\epsilon = 0$, the cumulative distribution functions of $X$ and $Y$ are identical, and thus the bound becomes exact, yielding $\mathrm{CVaR}_\alpha(X) = \mathrm{CVaR}_\alpha(Y)$. 
\begin{theorem}\label{thm:cvar_bound_convergence}
	Under the definitions of $X$, $Y$, $\epsilon$, and $\alpha$ specified in Theorem~\ref{thm:cvar_bound_v2}, the lower and upper bounds established therein converge to $\operatorname{CVaR}_\alpha(X)$ as $\epsilon \to 0$.
\end{theorem}
In the non-extreme cases, it holds that $\alpha>\epsilon$ and $\alpha+\epsilon\leq 1$, reducing the bounding problem to \eqref{eq:general_cvar_upper_bound_1} and \eqref{eq:general_cvar_lower_bound_1}. As we will see in the next section, $\epsilon$ could be very close to zero in practice. For $\alpha > \epsilon$, the lower bound on $\mathrm{CVaR}_\alpha(X)$ is a weighted average of the $\mathrm{CVaR}$ of $Y$ at a confidence level shifted by the distributional discrepancy $\epsilon$ and the maximum of the supports of $X$ and $Y$, where the weights are proportional to the amount of distributional discrepancy. For $\alpha+\epsilon \leq 1$, by restructuring \eqref{eq:general_cvar_lower_bound_1} as follows, 
\begin{equation}
	\begin{aligned}
		CVaR_{\alpha}(X) &\geq CVaR_{\alpha+\epsilon}(Y) \\
		&+ \frac{\epsilon}{\alpha}(CVaR_{\alpha+\epsilon}(Y)- CVaR_{\epsilon}(Y)).
	\end{aligned}
\end{equation}
The lower bound corresponds to the $\operatorname{CVaR}$ of $Y$ evaluated at a confidence level adjusted by the distributional discrepancy $\epsilon$, augmented by a correction term that is proportional to the distributional discrepancy.

Theorem~\ref{thm:cvar_bound_v2} assumes that the parameter $\epsilon$ provides an upper bound on the pointwise difference between the cumulative distribution functions $F_X$ and $F_Y$. As illustrated in Figure~\ref{fig:eps_bound}, this bound is particularly conservative in the vicinity of $x = -1$, where the actual discrepancy between $F_X$ and $F_Y$ is significantly smaller than the global bound $\epsilon$. Ideally, a tighter bound on $F_X(x)$ would allow for variation with respect to $x$, rather than relying on a uniform constant. Specifically, one seeks a pointwise bound of the form $F_X(x) \leq F_Y(x) + g(x)$ for some non-negative function $g$, as illustrated in Figure~\ref{fig:g_bound}.

\begin{figure*} % Use figure* to span both columns
	\centering
	\subfloat[]{\includegraphics[width=0.4\textwidth]{two_norm_cdfs_eps.png}\label{fig:eps_bound}}	\subfloat[]{\includegraphics[width=0.4\textwidth]{two_norm_cdfs.png}\label{fig:g_bound}}
	\caption{(a) and (b) are illustrations of the bounds on $F_X(x)$ from Theorems \ref{thm:cvar_bound_v2} and \ref{thm:tight_cvar_lower_bound}. Since $g$ from Theorem \ref{thm:tight_cvar_lower_bound} depends on $x$, the bound $F_Y(x)+g(x)$ in (b) is tighter than $F_Y(x)+\epsilon$ in (a).}
\end{figure*}


\begin{theorem}\label{thm:tight_cvar_lower_bound}
	(Tighter CVaR Lower Bound) Let $\alpha \in (0,1)$, $X$ and $Y$ be random variables. Define the a random variable $Y^L$ such that $F_{Y^L}(y) \triangleq min(1, F_Y(y) + g(y))$ for $g:\mathbb{R}\rightarrow [0, \infty)$. Assume 
	$lim_{x \rightarrow -\infty}g(x)=0$, $g$ is continuous from the right and monotonic increasing.
	If $\forall x\in \mathbb{R}, F_X(x)\leq F_Y(x)+g(x)$, then $F_{Y^L}$ is a CDF and $CVaR_\alpha(Y^L)\leq CVaR_\alpha(X).$
\end{theorem}

Theorem \ref{thm:tight_cvar_lower_bound} defines a random variable $Y^L$, constructed from $Y$ and the function $g$, such that the distributional discrepancy between $Y^L$ and $X$ is determined explicitly by $g$ rather than being uniformly bounded as in Theorem \ref{thm:cvar_bound_v2}. In addition, it offers a criterion for determining whether a given function $g$ can be used to derive a lower bound on the CVaR. This bound extends the lower bound established in Theorem \ref{thm:cvar_bound_v2}, which is obtained when $g(x)$ is constant and equal to $\epsilon$ for all $x \in \mathbb{R}$. 

Note that in Theorem \ref{thm:tight_cvar_lower_bound}, the function $g$ is assumed to be non-decreasing and right-continuous. If one assumes only that $|F_X(x) - F_Y(x)| \leq g(x)$ for some function $g : \mathbb{R} \to [0, \infty)$, which is not necessarily monotonic or continuous, then the most general form of the CVaR bounds is given by
\begin{equation}\label{eq:general_g_cvar_lower_bound}
	\begin{aligned}
		&CVaR_\alpha(X)=\frac{1}{\alpha}\int_{1-\alpha}^1 \inf\{z\in \mathbb{R}:F_X(z)\geq \tau\}d\tau \\
		&\geq \frac{1}{\alpha}\int_{1-\alpha}^1 \inf\{z\in \mathbb{R}:F_Y(z)+g(z)\geq \tau\}d\tau,
	\end{aligned}
\end{equation}
\begin{equation}\label{eq:general_g_cvar_upper_bound}
	\begin{aligned}
		&CVaR_\alpha(X)=\frac{1}{\alpha}\int_{1-\alpha}^1 \inf\{z\in \mathbb{R}:F_X(z)\geq \tau\}d\tau \\
		&\leq \frac{1}{\alpha}\int_{1-\alpha}^1 \inf\{z\in \mathbb{R}:F_Y(z)-g(z)\geq \tau\}d\tau.
	\end{aligned}
\end{equation}

Another option is to specify the distributional discrepancy through the density functions underlying the cumulative distribution functions. Let $f_x$ and $f_y$ be the probability density functions of $X$ and $Y$, respectively, and let $h : \mathbb{R} \to [0, \infty)$ describe the pointwise discrepancy between them. Theorem \ref{thm:lower_cvar_bound_using_density_bound} specifies conditions on the function $h$ under which a lower bound for the CVaR of $X$ can be obtained. A key advantage is that this bound takes the form of the CVaR of a random variable, enabling its estimation with performance guarantees via CVaR concentration bounds given in  theorems \ref{thm:brown_bounds}, \ref{thm:thomas_lower_bound}, \ref{thm:thomas_upper_bound} and Theorem \ref{thm:ecdf_cvar_bound}.

\begin{theorem}\label{thm:lower_cvar_bound_using_density_bound}
	Let $\alpha \in (0,1)$, $X$ and $Y$ random variables. Define $h:\mathbb{R}\rightarrow [0,\infty)$ to be a continuous function, $g(z)\triangleq\int_{-\infty}^z h(x)dx$ and $Y^L$ to be a random variable such that $F_{Y^L}(y)\triangleq\min(1, F_{Y}(y) + g(y))$. If $\lim_{z\rightarrow -\infty} g(z)=0$ and $\forall x\in \mathbb{R},f_x(x)\leq f_y(z) + h(x)$, then $F_{Y^L}$ is a CDF and $CVaR_\alpha(Y^L)\leq CVaR_\alpha(X)$.
\end{theorem}

Theorem \ref{thm:lower_cvar_bound_using_density_bound} is obtained as a corollary of Theorem \ref{thm:tight_cvar_lower_bound} by defining a function $g$, as illustrated in Figure \ref{fig:g_bound}, in terms of the function $h$. This construction demonstrates that the function $g$ in Theorem \ref{thm:tight_cvar_lower_bound} can be generated from a broad class of density discrepancy functions.

\section{Concentration Inequalities}\label{sec:con_ineq}
\addtocounter{section}{1}
In this section, we derive concentration inequalities for $\operatorname{CVaR}_\alpha(X)$ based on samples drawn from an auxiliary random variable $Y$. A notable special case of these inequalities arises when $Y$ is taken to follow the empirical cumulative distribution function (ECDF) of $X$. In such cases, concentration inequalities for $\operatorname{CVaR}_\alpha(X)$ are obtained.

Let $X_1, \ldots, X_n \overset{i.i.d.}{\sim} F_X$, where $F_X$ is the CDF of a random variable $X$. The ECDF based on these samples is defined by
\begin{equation}
	\hat{F}_X(x)=\frac{1}{n} \sum_{i=1}^n 1_{X_i\leq x},
\end{equation}
for $x\in \mathbb{R}$. Let $E_{\hat{F}_X}[X]$ denote the expectation with respect to $\hat{F}_X$, and let $C_\alpha^{\hat{F}_X}$ denote the CVaR computed under the empirical distribution $\hat{F}_X$.

%\begin{theorem}
%	\label{thm:ecdf_cvar_bound}
%	\KP{[Thomas assumes only one side bound.]}Let $X\in [a,b]$ be some random variable for $a,b\in \mathbb{R}$, and $\alpha\in (0, 1],\delta \in (0, 1), \epsilon=\sqrt{\ln(1/\delta)/(2n)}$. Let $X_1, \dots, X_n \overset{iid}{\sim} F_X$ be random variables that define the ECDF $\hat{F}_X$. 
%	\begin{enumerate}
%		\item If $\alpha > \epsilon$, then $P(CVaR_\alpha(X) \leq (1-\frac{\epsilon}{\alpha})C_{\alpha-\epsilon}^{\hat{F}_X} + \frac{\epsilon}{\alpha}b)> 1-\delta$.
%		
%		\item If $\alpha\leq \epsilon$ then $CVaR_\alpha(X)\leq b$
%		
%		\item If $\alpha+\epsilon < 1$, then $P(CVaR_\alpha(X)\geq (1+\frac{\epsilon}{\alpha})C_{\alpha+\epsilon}^{\hat{F}_X}-\frac{\epsilon}{\alpha}C_{\epsilon}^{\hat{F}_X})> 1-\delta$.
%		
%		\item If $\alpha+\epsilon \geq 1$, then $P(CVaR_\alpha(X) \geq \frac{1}{\alpha}[(\alpha+\epsilon - 1)a + E_{\hat{F}_X}[X] - \epsilon C_\epsilon^{\hat{F}_X}]) > 1-\delta$.
%	\end{enumerate}
%\end{theorem}
\begin{theorem}
	\label{thm:ecdf_cvar_bound}
	Let $X$ a random variable, $\alpha\in (0, 1],\delta \in (0, 1), \epsilon=\sqrt{\ln(1/\delta)/(2n)}$. Let $X_1, \dots, X_n \overset{iid}{\sim} F_X$ be random variables that define the ECDF $\hat{F}_X$. 
	\begin{enumerate}
		\item \textbf{Upper Bound:} If $P(X\leq b)=1$, then 
			\begin{enumerate}
				\item If $\alpha > \epsilon$ then $P(CVaR_\alpha(X) \leq (1-\frac{\epsilon}{\alpha})C_{\alpha-\epsilon}^{\hat{F}_X} + \frac{\epsilon}{\alpha}b)> 1-\delta$.
				
				\item If $\alpha\leq \epsilon$ then $CVaR_\alpha(X)\leq b$
			\end{enumerate}
			
		\item \textbf{Lower Bound:} If $P(X\geq a)=1$, then 
			\begin{enumerate}
				\item If $\alpha+\epsilon < 1$, then $P(CVaR_\alpha(X)\geq (1+\frac{\epsilon}{\alpha})C_{\alpha+\epsilon}^{\hat{F}_X}-\frac{\epsilon}{\alpha}C_{\epsilon}^{\hat{F}_X})> 1-\delta$.
				
				\item If $\alpha+\epsilon \geq 1$, then $P(CVaR_\alpha(X) \geq \frac{1}{\alpha}[(\alpha+\epsilon - 1)a + E_{\hat{F}_X}[X] - \epsilon C_\epsilon^{\hat{F}_X}]) > 1-\delta$.
			\end{enumerate}
	\end{enumerate}
\end{theorem}
Theorem \ref{thm:ecdf_cvar_bound} provides concentration inequalities for $\operatorname{CVaR}_\alpha(X)$ in a form that enables the user to specify a desired bound consistency level $\delta$, which determines the probability that the bound holds. This result follows as a corollary of Theorem \ref{thm:cvar_bound_v2}, in which the auxiliary random variable $Y$ is instantiated as the ECDF of $X$, whereas $X$ denotes the underlying true random variable, which is inaccessible in practice. The distributional discrepancy required by Theorem \ref{thm:cvar_bound_v2}, denoted by $\epsilon$, is controlled in Theorem \ref{thm:ecdf_cvar_bound} via the Dvoretzky–Kiefer–Wolfowitz (DKW) inequality \cite{DKW_inequality}. The DKW inequality ensures that the supremum distance between the true CDF and the ECDF converges to zero at a rate of order $1/\sqrt{n}$ as the number of samples $n$ increases. Corollary \ref{thm:x_concentration_bound_asymptotic_convergence} establishes the asymptotic convergence of the bounds given in Theorem \ref{thm:ecdf_cvar_bound} as the sample size tends to infinity.
%\begin{corollary}\label{thm:x_concentration_bound_asymptotic_convergence}
%	\KP{[adjust proof]}\KP{[Thomas assumes only one side bound.]}Let $X\in [a,b]$ be some random variable for $a,b\in \mathbb{R}$, and $\alpha\in (0, 1],\delta \in (0, 1), \epsilon=\sqrt{\ln(1/\delta)/(2n)}$. Let $X_1, \dots, X_n \overset{iid}{\sim} F_X$ be random variables that define the ECDF $\hat{F}_X$. Denote by $U(n)$ and $L(n)$ the upper and lower bounds respectively from Theorem \ref{thm:ecdf_cvar_bound}, where $n$ is the number of samples, then 
%	\begin{equation}
%		\lim_{n\rightarrow \infty} L(n)=CVaR_\alpha(X), \lim_{n\rightarrow \infty} U(n)=CVaR_\alpha(X)
%	\end{equation}
%\end{corollary}
\begin{corollary}\label{thm:x_concentration_bound_asymptotic_convergence}
	Let $X$ be a random variable, $\alpha\in (0, 1],\delta \in (0, 1), \epsilon=\sqrt{\ln(1/\delta)/(2n)}, a\in \mathbb{R}, b\in \mathbb{R}, \eta>0$. Let $X_1, \dots, X_n \overset{iid}{\sim} F_X$ be random variables that define the ECDF $\hat{F}_X$. Denote by $U(n)$ and $L(n)$ the upper and lower bounds respectively from Theorem \ref{thm:ecdf_cvar_bound}, where $n$ is the number of samples, then 
	\begin{enumerate}
		\item If $P(X\leq b)=1$, then $\lim_{n\rightarrow \infty} U(n)\overset{a.s}{=}CVaR_\alpha(X)$.
		\item If $P(X\geq a)=1$, then $\lim_{n\rightarrow \infty} L(n)\overset{a.s}{=}CVaR_\alpha(X)$.
	\end{enumerate}
	where a.s denotes almost sure convergence.
\end{corollary}
%\begin{proof}
%	\begin{equation}
%		\begin{aligned}
%			&\lim_{n \to \infty} U(n)\overset{a.s}{=} \lim_{n \to \infty} (1-\frac{\epsilon}{\alpha})C_{\alpha-\epsilon}^{\hat{F}_X} + \frac{\epsilon}{\alpha}b \\
%			&\overset{a.s}{=}\lim_{n \to \infty} C_{\alpha-\epsilon}^{\hat{F}_X}=\lim_{n \to \infty} C_{\alpha}^{\hat{F}_X} = C_{\alpha}^{F_X}
%		\end{aligned}
%	\end{equation}
%	where the first and second equalities follow from the fact that $\epsilon \to 0$ as $n \to \infty$; the third equality holds by the continuity of CVaR with respect to $\alpha$; and the fourth equality holds by the almost sure convergence of the empirical CVaR of i.i.d. samples to the true CVaR. For the same reason,
%	\begin{equation}
%		\begin{aligned}
%			&\lim_{n\rightarrow \infty} L(n)\overset{a.s}{=} \lim_{n\rightarrow \infty} (1+\frac{\epsilon}{\alpha})C_{\alpha+\epsilon}^{\hat{F}_X}-\frac{\epsilon}{\alpha}C_{\epsilon}^{\hat{F}_X} \\
%			&\overset{a.s}{=} \lim_{n\rightarrow \infty} C_{\alpha+\epsilon}^{\hat{F}_X}
%			\overset{a.s}{=} \lim_{n\rightarrow \infty} C_{\alpha}^{\hat{F}_X}
%			\overset{a.s}{=} \lim_{n\rightarrow \infty} C_{\alpha}^{F_X}.
%		\end{aligned}
%	\end{equation}
%%	Let $Y \sim F_Y=\hat{F}_X$ be an auxiliary random variable, where $\hat{F}_X$ is the ECDF of $X$ defined by $X_1,\dots, X_n$. If $\sup_{x \in \mathbb{R}} \big( F_Y(x) - F_X(x) \big) 
%%	\leq \epsilon$ then the conditions of Theorem \ref{thm:cvar_bound_v2} are satisfied, and we get the exact bounds of Theorem \ref{thm:ecdf_cvar_bound}. As $n \to \infty$, we have $\epsilon \to 0$, and by Theorem~\ref{thm:cvar_bound_convergence}, it follows that $\lim_{n \to \infty} U(n) = \operatorname{CVaR}_\alpha(X)$. Likewise, if $\sup{x \in \mathbb{R}} \big(F_Y(x) - F(x)\big) \to 0$, then Theorem~\ref{thm:cvar_bound_convergence} implies that $\lim_{n \to \infty} L(n) = \operatorname{CVaR}_\alpha(X)$.
%%	
%%	From DKW \cite{DKW_inequality} inequality the following inequalities can be derived \cite{pmlr-v97-thomas19a}
%%	\begin{equation}\label{eq:helper_2}
%%		\Pr \left( \sup_{x \in \mathbb{R}} \big( F_Y(x) - F_X(x) \big) 
%%		\leq \epsilon \right) \geq 1 - \delta,
%%	\end{equation}
%%	\begin{equation}
%%		\Pr \left( \sup_{x \in \mathbb{R}} \big( F_Y(x) - F_X(x) \big) 
%%		\geq \epsilon \right) \geq 1 - \delta.
%%	\end{equation}  
%%	and therefore 
%%	\begin{equation}
%%		\begin{aligned}
%%			&P(|U(n)-CVaR_\alpha(X)|>\eta)\\
%%			&=P(|U(n)-CVaR_\alpha(X)|>\eta|\sup_{x \in \mathbb{R}} \big( F_Y(x) - F_X(x) \big) 
%%			\leq \epsilon) \\
%%			&\times P(\sup_{x \in \mathbb{R}} \big( F_Y(x) - F_X(x) \big) 
%%			\leq \epsilon) \\
%%			&+P(|U(n)-CVaR_\alpha(X)|>\eta|\sup_{x \in \mathbb{R}} \big( F_Y(x) - F_X(x) \big) 
%%			> \epsilon) \\
%%			&\times P(\sup_{x \in \mathbb{R}} \big( F_Y(x) - F_X(x) \big) 
%%			> \epsilon) \\
%%			\geq 
%%		\end{aligned}
%%	\end{equation}
%\end{proof}


%\begin{proof}
%	\KP{[rephrase]}Theorem \ref{thm:cvar_bound_convergence} shows that as the distributional discrepancy $\epsilon \to 0$ \KP{[under the setting of Theorem \ref{thm:cvar_bound_v2}]}, we have $\mathrm{CVaR}\alpha(X) = \mathrm{CVaR}\alpha(Y)$. Furthermore, Theorem \ref{thm:ecdf_cvar_bound} can be regarded as a special case of Theorem \ref{thm:cvar_bound_v2}, where the random variable $Y$ corresponds to the ECDF of $X$. Since the distributional discrepancy $\epsilon$ in Theorem \ref{thm:ecdf_cvar_bound} converges to zero as a function of the sample size $n$, the claimed result follows directly from Theorem \ref{thm:cvar_bound_convergence}.
%\end{proof}

The concentration bounds established by \cite{pmlr-v97-thomas19a} coincide with those given in Theorem \ref{thm:ecdf_cvar_bound}, rendering the results of \cite{pmlr-v97-thomas19a} a special case of Theorem \ref{thm:ecdf_cvar_bound}. This equivalence arises because both Theorem \ref{thm:ecdf_cvar_bound} and \cite{pmlr-v97-thomas19a} derive concentration bounds for $\operatorname{CVaR}_\alpha(X)$ by constructing an alternative CDF that stochastically dominates the true distribution $F_X$, employing the DKW inequality to control the discrepancy. The principal distinction between the two results lies in the formulation of the $\operatorname{CVaR}$ bound: \cite{pmlr-v97-thomas19a} express the bound through a sum of reweighted order statistics (theorems \ref{thm:thomas_upper_bound} and \ref{thm:thomas_lower_bound}), resulting in a more intricate form, whereas Theorem \ref{thm:ecdf_cvar_bound} presents a more interpretable bound in terms of $\operatorname{CVaR}$. The interpretability of these bounds constitutes a contribution of this paper.

Theorem \ref{thm:x_concentration_bound_asymptotic_convergence} provides concentration inequalities for the theoretical $CVaR_\alpha(X)$ based on samples drawn from an auxiliary random variable $Y$, assuming only a bound on the distributional discrepancy between $X$ and $Y$.
%\begin{theorem}
%	\label{thm:ecdf_y_bound_cvar_x}
%	Let $X$ and $Y$ be random variables with supports $X \in [a_X, b_X]$ and $Y \in [a_Y, b_Y]$, where $a_X, a_Y, b_X, b_Y \in \mathbb{R}$. Define $a = \max(a_X, a_Y)$, $b = \min(b_X, b_Y), \epsilon \in [0, 1]$, and $\eta = \sqrt{\ln(1/\delta)/(2n)}$. Let $\epsilon'=min(\epsilon+\eta, 1)$. Let $Y_1, \dots, Y_n$ be independent and identically distributed samples from $F_Y$, and denote by $\hat{F}_Y$ the associated empirical cumulative distribution function.
%	\begin{enumerate}
%		\item If $\forall z\in \mathbb{R},F_Y(z)-F_X(z)\leq \epsilon$ and $\alpha > \epsilon'$ then 
%		\begin{equation}\label{eq:general_cvar_upper_bound_1}
%			\begin{aligned}
%				P\Big(CVaR_\alpha(X) &\leq \frac{\epsilon'}{\alpha}max(b_X,b_Y) \\
%				&+ (1 - \frac{\epsilon'}{\alpha})CVaR_{\alpha-\epsilon'}^{\hat{F}_Y^L}\Big)>1-\delta.
%			\end{aligned}
%		\end{equation}
%		
%		\item If $\forall z\in \mathbb{R},F_Y(z)-F_X(z)\leq \epsilon$ and $\alpha \leq \epsilon'$ then $CVaR_\alpha(X)\leq max(b_X,b_Y)$.
%		
%		\item If $\forall z\in \mathbb{R},F_X(z)-F_Y(z)\leq \epsilon$ and $\alpha+\epsilon' \leq 1$ then 
%		\begin{equation}\label{eq:general_cvar_lower_bound_1}
%			\begin{aligned}
%				P\Big(CVaR_{\alpha}(X)&\geq (1+\frac{\epsilon'}{\alpha})CVaR_{\alpha+\epsilon'}^{\hat{F}_Y^L} \\
%				&-\frac{\epsilon'}{\alpha} CVaR_{\epsilon'}^{\hat{F}_Y^L}\Big)>1-\delta
%			\end{aligned}
%		\end{equation}
%		
%		\item If $\forall z\in \mathbb{R},F_X(z)-F_Y(z)\leq \epsilon$ and $\alpha+\epsilon' > 1$ then 
%		\begin{equation}
%			\begin{aligned}
%				P\Big(CVaR_\alpha(X) &\geq \mathbb{E}_{\hat{F}_Y^L}[Y] - \epsilon' CVaR_{\epsilon'}^{\hat{F}_Y} \\
%				&+ (\alpha+\epsilon'-1)a_{min} \Big) > 1-\delta.
%			\end{aligned}
%		\end{equation}
%	\end{enumerate}
%\end{theorem}
\begin{theorem}
	\label{thm:ecdf_y_bound_cvar_x}
	Let $X$ and $Y$ be random variables, $\epsilon \in [0, 1]$, and $\eta = \sqrt{\ln(1/\delta)/(2n)}, \epsilon'=min(\epsilon+\eta, 1)$. Let $Y_1, \dots, Y_n$ be independent and identically distributed samples from $F_Y$, and denote by $\hat{F}_Y$ the associated empirical cumulative distribution function.
	\begin{enumerate}
		\item \textbf{Upper Bound:} If $\forall z\in \mathbb{R},F_Y(z)-F_X(z)\leq \epsilon$ and $P(X\leq b_X)=1, P(Y\leq b_Y)=1$, then
		
		\begin{enumerate}
			\item If $\alpha > \epsilon'$ then \begin{equation}\label{eq:general_cvar_upper_bound_1}
				\begin{aligned}
					P\Big(CVaR_\alpha(X) &\leq \frac{\epsilon'}{\alpha}max(b_X,b_Y) \\
					&+ (1 - \frac{\epsilon'}{\alpha})CVaR_{\alpha-\epsilon'}^{\hat{F}_Y^L}\Big)>1-\delta.
				\end{aligned}
			\end{equation}
			
			\item If $\alpha \leq \epsilon'$, then $CVaR_\alpha(X)\leq max(b_X,b_Y)$
		\end{enumerate}
		
		\item \textbf{Lower Bound:} If $\forall z\in \mathbb{R},F_X(z)-F_Y(z)\leq \epsilon$ and $P(X\geq a_X)=1,P(Y\geq a_Y)=1$, then \begin{enumerate}
			\item If $\alpha+\epsilon' \leq 1$, then \begin{equation}\label{eq:general_cvar_lower_bound_1}
				\begin{aligned}
					P\Big(CVaR_{\alpha}(X)&\geq (1+\frac{\epsilon'}{\alpha})CVaR_{\alpha+\epsilon'}^{\hat{F}_Y^L} \\
					&-\frac{\epsilon'}{\alpha} CVaR_{\epsilon'}^{\hat{F}_Y^L}\Big)>1-\delta.
				\end{aligned}
			\end{equation}
			
			\item If $\alpha+\epsilon' > 1$, then \begin{equation}
				\begin{aligned}
					P\Big(CVaR_\alpha(X) &\geq \mathbb{E}_{\hat{F}_Y^L}[Y] - \epsilon' CVaR_{\epsilon'}^{\hat{F}_Y} \\
					&+ (\alpha+\epsilon'-1)a_{min} \Big) > 1-\delta.
				\end{aligned}
			\end{equation}
		\end{enumerate}
	\end{enumerate}
\end{theorem}
\begin{proof}
	We begin by establishing that $\epsilon'$ bounds the distributional discrepancy between $\hat{F}_Y^L$ and $F_X$, under the assumption that $\sup_x \big( \hat{F}_Y^L(x) - F_Y(x) \big) \leq \epsilon$. Let $x\in \mathbb{R}$,
	\begin{equation}
		\begin{aligned}
			&\hat{F}_Y^L(x)-F_X(x)=\hat{F}_Y^L(x)-F_Y(x) + F_Y(x)-F_X(x) \\
			&\leq |\hat{F}_Y^L(x)-F_Y(x)| + |F_Y(x)-F_X(x)| \leq \epsilon + \eta=\epsilon'.
		\end{aligned}
	\end{equation}
	Assume that $\alpha > \epsilon'$. Note that, conditional on the event $\sup_{z \in \mathbb{R}} \big( \hat{F}_Y(z) - F_X(z) \big) \leq \epsilon'$,
	the upper bound in Theorem~\ref{thm:cvar_bound_v2} holds deterministically. Consequently, the probability that
	\begin{equation}
		CVaR_\alpha(X) \leq \frac{\epsilon'}{\alpha}max(b_X,b_Y) 
		+ (1 - \frac{\epsilon'}{\alpha})CVaR_{\alpha-\epsilon'}^{\hat{F}_Y^L}
	\end{equation}
	holds is equal to one. From the law of total probability, 
	\begin{equation}
		\begin{aligned}
			&P\Big(CVaR_\alpha(X) \leq \frac{\epsilon'}{\alpha}max(b_X,b_Y) + (1 - \frac{\epsilon'}{\alpha})CVaR_{\alpha-\epsilon'}^{\hat{F}_Y^L}\Big) \\
			&=P\Big(CVaR_\alpha(X) \leq \frac{\epsilon'}{\alpha}max(b_X,b_Y) + (1 - \frac{\epsilon'}{\alpha})CVaR_{\alpha-\epsilon'}^{\hat{F}_Y^L} \\
			&\Big|\sup_{z \in \mathbb{R}} \big( \hat{F}_Y(z) - F_X(z) \big) \leq \epsilon'\Big) 
			P(\sup_{z \in \mathbb{R}} \big( \hat{F}_Y(z) - F_X(z) \big) \leq \epsilon') \\
			&+P\Big(CVaR_\alpha(X) \leq \frac{\epsilon'}{\alpha}max(b_X,b_Y) + (1 - \frac{\epsilon'}{\alpha})CVaR_{\alpha-\epsilon'}^{\hat{F}_Y^L} \\
			&\Big|\sup_{z \in \mathbb{R}} \big( \hat{F}_Y(z) - F_X(z) \big) > \epsilon'\Big) 
			P(\sup_{z \in \mathbb{R}} \big( \hat{F}_Y(z) - F_X(z) \big) > \epsilon')\\
			&=P(\sup_{z \in \mathbb{R}} \big( \hat{F}_Y(z) - F_X(z) \big) \leq \epsilon') \\
			&+P\Big(CVaR_\alpha(X) \leq \frac{\epsilon'}{\alpha}max(b_X,b_Y) + (1 - \frac{\epsilon'}{\alpha})CVaR_{\alpha-\epsilon'}^{\hat{F}_Y^L} \\
			&\Big|\sup_{z \in \mathbb{R}} \big( \hat{F}_Y(z) - F_X(z) \big) > \epsilon'\Big) 
			P(\sup_{z \in \mathbb{R}} \big( \hat{F}_Y(z) - F_X(z) \big) > \epsilon') \\
			&\geq P(\sup_{z \in \mathbb{R}} \big( \hat{F}_Y(z) - F_X(z) \big) \leq \epsilon')
		\end{aligned}
	\end{equation}
	From DKW \cite{DKW_inequality} inequality the following inequalities can be derived \cite{pmlr-v97-thomas19a}
	\begin{equation}\label{eq:helper_2}
		\Pr \left( \sup_{x \in \mathbb{R}} \big( \hat{F}(x) - F(x) \big) 
		\leq \sqrt{\frac{\ln(1/\delta)}{2n}} \right) \geq 1 - \delta,
	\end{equation}
	\begin{equation}
		\Pr \left( \sup_{x \in \mathbb{R}} \big( \hat{F}(x) - F(x) \big) 
		\geq \sqrt{\frac{\ln(1/\delta)}{2n}} \right) \geq 1 - \delta.
	\end{equation}
	\begin{equation}
		\begin{aligned}
			&P(\sup_{z \in \mathbb{R}} \big( \hat{F}_Y(z) - F_X(z) \big) \leq \epsilon') \\
			&\geq P(\sup_{z \in \mathbb{R}} \big( \hat{F}_Y(z) - F_Y(z) \big) + \sup_{z \in \mathbb{R}} \big( F_Y(z) - F_X(z) \big)\leq \epsilon') \\
			&\geq P(\sup_{z \in \mathbb{R}} \big( \hat{F}_Y(z) - F_Y(z) \big) + \epsilon \leq \epsilon + \eta) \\
			&=P(\sup_{z \in \mathbb{R}} \big( \hat{F}_Y(z) - F_Y(z) \big)\leq \eta) > 1-\delta.
		\end{aligned}
	\end{equation}
	The first inequality follows from the triangle inequality; the second holds since $\epsilon$ bounds the distributional discrepancy between $X$ and $Y$; and the third follows from \eqref{eq:helper_2}.
	
	As with the preceding equations, all bounds in Theorem~\ref{thm:cvar_bound_v2} hold deterministically for a given distributional discrepancy, where the discrepancy between $\hat{F}_Y$ and $F_X$ is $\epsilon'$. Consequently, the remaining probabilistic guarantees hold with probability at least $1 - \delta$.
\end{proof}

The parameter $\epsilon$ in Theorem \ref{thm:x_concentration_bound_asymptotic_convergence} captures the distributional discrepancy between $X$ and $Y$, consistent with its role in Theorem \ref{thm:cvar_bound_v2}. Additionally, the theorem introduces $\epsilon'$ to represent the discrepancy between $F_X$ and $\hat{F}_Y$, where $\hat{F}_Y$ denotes the empirical CDF of $Y$ constructed from the sample $\{Y_i\}_{i=1}^n$. The parameter $\eta$ accounts for the additional distributional discrepancy beyond $\epsilon$, and captures the estimation error in approximating $F_Y$ using the empirical sample $\{Y_i\}_{i=1}^n$. By combining Theorem \ref{thm:cvar_bound_v2} with the DKW inequality \cite{DKW_inequality}, we obtain probabilistic guarantees for the estimated bounds.

%\begin{figure}[t]
%	\centering
%	\includegraphics[width=0.4\textwidth]{bounded_actions_plot.png}
%	\caption{
%		Illustration of risk-averse action selection via CVaR bounds.
%	}
%	\label{fig:bounded_actions_plot}
%\end{figure}

\section{Computationally Efficient Sampling}
In situations where the CVaR of a random variable $X$ is computationally intractable, the bounds derived in the preceding sections may serve as a practical alternative for efficient estimation. Specifically, one may select an auxiliary random variable $Y$ that is more amenable to sampling, and employ the established bounds to approximate the CVaR of $X$ using samples drawn from $Y$. This approach is particularly relevant in settings such as online planning, where computational efficiency is critical \cite{Lev-Yehudi_Barenboim_Indelman_2024, ijcai2022p637}.

%Online decision-making for agents often requires recurrent sampling to evaluate the quality of actions, a process that can be computationally intensive. In the context of Partially Observable Markov Decision Processes (POMDPs), state-of-the-art planning algorithms simulate prospective trajectories up to a fixed computational budget, incrementally expanding a decision tree as time permits \cite{silver2010monte, sunberg2018online}. At each decision node, the agent samples from the state-transition and observation models—which may be represented by complex neural networks—to estimate the CVaR associated with each available action. A strategy to reduce this burden is to compute bounds on key decision metrics rather than estimating them precisely. 
%
%To enhance computational efficiency, the original models may be replaced with surrogate dynamics that are more tractable for sampling. For instance, \cite{levyehudi2023simplifying} perform planning directly with respect to a surrogate model, while \cite{barenboim2022adaptive} utilize bounds on the true value function derived from surrogate-based approximations, both relying on an expectation-based formulation. The use of such surrogate models necessitates a rigorous theoretical framework that ensures the original value function can be reliably bounded in terms of the surrogate-based one. The theoretical results presented in this paper establish such a foundation by deriving conditions under which CVaR bounds computed from surrogate models can safely inform action selection. These bounds allow the agent to estimate confidence intervals on CVaR values; when the intervals corresponding to different actions are disjoint, the agent may safely prefer the more risk-averse action without the need to compute the exact CVaR.
%
%Figure~\ref{fig:bounded_actions_plot} illustrates two representative scenarios. In the left panel (Non-overlapping Bounds), the confidence intervals for the CVaR estimates of the two actions do not intersect, allowing the agent to decisively prefer Action 1. In contrast, the right panel (Overlapping Bounds) shows a case where the intervals overlap, preventing a clear comparison. In such cases, further evaluation using the full models or additional samples is required to inform the decision.

\section{Experiments}
In this simulation study, we demonstrate how the CVaR of an inaccessible or computationally expensive distribution can be bounded using an auxiliary distribution that is more tractable. We then verify that the concentration bounds for $CVaR_\alpha(X)$ established in Theorem~\ref{thm:ecdf_cvar_bound} coincide with those of \cite{pmlr-v97-thomas19a}, which arise as a special case of our more general results. In our experiments, both approaches exhibit identical empirical performance.
%\begin{figure*}
%	\centering
%	\subfloat[]{\includegraphics[width=0.4\textwidth]{cvar_convergence_with_bounds.png}\label{fig:cvar_bounds_convergence}}	\subfloat[]{\includegraphics[width=0.4\textwidth]{sampling_time_comparison.png}\label{fig:cvar_bounds_running_time}}
%	\caption{\KP{[add]}}
%\end{figure*}
\begin{figure*}[t]
	\centering
	\subfloat[]{\includegraphics[width=0.3\textwidth]{cvar_convergence_with_bounds.png}\label{fig:cvar_bounds_convergence}}
	\subfloat[]{\includegraphics[width=0.3\textwidth]{sampling_time_comparison.png}\label{fig:sampling_running_time}}
	\subfloat[]{\includegraphics[width=0.3\textwidth]{cvar_computation_time_comparison.png}\label{fig:cvar_time_comparison_1}}
	\caption{Figure \ref{fig:cvar_bounds_convergence} shows the bounds established in Theorem \ref{thm:ecdf_y_bound_cvar_x}, demonstrating how a simple truncated Normal distribution can be employed to bound the CVaR of a truncated GMM. Figure \ref{fig:sampling_running_time} presents a comparison of the sampling times for each distribution, while Figure \ref{fig:cvar_time_comparison_1} reports the total computation time required to sample from the truncated GMM and estimate its CVaR, in contrast to sampling from the truncated Normal surrogate and computing the associated concentration bounds. All of the experiments were configured with confidence level $\alpha=0.2$, and probability of error $\delta=0.05$. Both Figures \ref{fig:sampling_running_time} and \ref{fig:cvar_time_comparison_1} display 95\% confidence intervals around the empirical means.}
\end{figure*}

Specifically, we compare a truncated Gaussian Mixture Model (GMM) and its corresponding truncated Normal approximation for CVaR estimation under bounded support. Such a setting arises, for example, in POMDP planning, where a computationally expensive model used by the agent during online planning is replaced with a more tractable surrogate model, thereby improving the agent’s decision-making speed \cite{Lev-Yehudi_Barenboim_Indelman_2024}. In these settings, the agent’s decisions may rely on bounds for the original value function that are derived from the tractable surrogate model \cite{ijcai2022p637}.

The GMM consists of five components with means $\mu_1 = 0.2$, $\mu_2 = -0.2$, $\mu_3 = -0.5$, $\mu_4 = 0.5$, $\mu_5 = 0.0$, variances $\sigma_1^2 = 0.5$, $\sigma_2^2 = 0.2$, $\sigma_3^2 = 0.1$, $\sigma_4^2 = 0.1$, $\sigma_5^2 = 0.3$, and weights $w_1 = 0.6$, $w_2 = 0.4$, $w_3 = 0.1$, $w_4 = 0.1$, $w_5 = 0.8$. The Normal approximation matches the GMM's mean and variance but cannot reproduce its multi-modal structure or outlier effects. All samples are truncated to the interval $[-1, 1]$, which reshapes the tails and slightly distorts boundary components. We compute CVaR at the 20\% quantile for sample sizes ranging from 100 to 10{,}000, with 100 independent repetitions per setting to ensure statistical reliability. The distributional discrepancy ($\epsilon$ in Theorem \ref{thm:ecdf_y_bound_cvar_x}) between the truncated GMM and the truncated Normal distribution is assessed via simulation, by estimating their respective cumulative distribution functions over a common set of bins and computing the maximum difference across all bins. This setup enables a direct assessment of the trade-off between computational efficiency and statistical accuracy when approximating a complex truncated mixture by a single truncated Normal distribution.

Figure \ref{fig:sampling_running_time} presents the sampling time comparison between the truncated GMM and the truncated Normal distribution, with an observed average time ratio of approximately 3.6 in favor of the Normal distribution. Figure \ref{fig:cvar_time_comparison_1} reports a total computational speedup of approximately 3.17 when comparing the process of sampling from the truncated GMM and estimating its CVaR to that of sampling from the truncated Normal and computing both upper and lower bounds as given in Theorem \ref{thm:ecdf_y_bound_cvar_x}. Figure \ref{fig:cvar_bounds_convergence} illustrates the convergence behavior of the CVaR bounds as a function of the number of samples, comparing estimates obtained from the truncated GMM with bounds derived from the truncated Normal surrogate. It is important to note that these bounds are obtained without requiring full knowledge of the underlying GMM distribution. Instead, they rely solely on the discrepancy between the corresponding CDFs. 

We evaluated the concentration bounds established by \cite{pmlr-v97-thomas19a} alongside our proposed bounds from Theorem \ref{thm:ecdf_cvar_bound} on a set of probability distributions: $\operatorname{Beta}(2, 2)$, $\operatorname{Beta}(0.5, 0.5)$, $\operatorname{Beta}(2, 5)$, $\operatorname{Beta}(5, 2)$, $\operatorname{Beta}(10, 2)$, $\operatorname{Beta}(2, 10)$, and the Laplace$(0, 1)$ distribution. These distributions were selected to match those used in the original study by \cite{pmlr-v97-thomas19a}, enabling a direct comparison under identical conditions. For each distribution, our bounds precisely coincide with those reported by \cite{pmlr-v97-thomas19a}, resulting in complete overlap between the two sets of bounds. The graphs exhibiting these results are available in the supplemental material.

\section{Conclusions}
This work presented interpretable concentration inequalities for $CVaR_\alpha(X)$, expressed in terms of the CVaR of an auxiliary random variable $Y$. The proposed bounds were validated through simulation studies, effectively bounding the CVaR of a truncated GMM by employing a less complex truncated Normal distribution. A comprehensive theoretical framework was developed for bounding CVaR using both uniform and non-uniform bounds on the cumulative and density functions of $X$ and $Y$. Notably, the state-of-the-art bounds introduced by \cite{pmlr-v97-thomas19a} emerge as a special case of the proposed general framework, which also offers a novel interpretation of their result. 

\bibliography{aaai2026}

\section{Supplemental Material}

\textbf{Note:} The theorems in this supplementary material have been reindexed; therefore, references to theorem numbers pertain only to the numbering used within this section. Consequently, the same theorem may be assigned a different number here than in the main paper.

\section{Proofs}
\addtocounter{section}{1}
\begin{theorem}\label{thm:cvar_bound_v2}
	Let $X$ and $Y$ be random variables and $\epsilon \in [0, 1]$. 
	\begin{enumerate}
		\item \textbf{Upper bound:} assume that $P(X\leq b_X)=1, P(Y\leq b_Y)=1$ and $\forall z\in \mathbb{R},F_Y(z)-F_X(z)\leq \epsilon$, then 
		\begin{enumerate}
			\item If $\alpha > \epsilon$, \begin{equation}\label{eq:general_cvar_upper_bound_1}
				CVaR_\alpha(X) \leq \frac{\epsilon}{\alpha}max(b_X,b_Y) + (1 - \frac{\epsilon}{\alpha})CVaR_{\alpha-\epsilon}(Y).
			\end{equation}
			
			\item If $\alpha \leq \epsilon$, then $CVaR_\alpha(X)\leq max(b_X,b_Y)$.
		\end{enumerate}
		
		\item \textbf{Lower Bound:} If $P(X\geq a_X)=1, P(Y\geq a_Y)=1$ and $\forall z\in \mathbb{R},F_X(z)-F_Y(z)\leq \epsilon$, then
		\begin{enumerate}
			\item If $\alpha+\epsilon \leq 1$, \begin{equation}\label{eq:general_cvar_lower_bound_1}
				CVaR_{\alpha}(X)\geq (1+\frac{\epsilon}{\alpha})CVaR_{\alpha+\epsilon}(Y)-\frac{\epsilon}{\alpha} CVaR_{\epsilon}(Y)
			\end{equation}
			
			\item If $\alpha+\epsilon > 1$ and $a_{min} = min(a_X, a_Y)$, \begin{equation}
				CVaR_\alpha(X) \geq \mathbb{E}[Y] - \epsilon CVaR_\epsilon(Y) + (\alpha+\epsilon-1)a_{min}
			\end{equation}
		\end{enumerate}
	\end{enumerate}
\end{theorem}
\begin{proof}
	The strategy of the proof is to construct two distributions derived from $F_Y$, denoted $F_Y^L$ and $F_Y^U$, such that $F_X$ is stochastically bounded between them; that is, $F_Y^L \leq F_X \leq F_Y^U$. Consequently, since CVaR is a coherent risk measure, it follows that
	\begin{equation}
		CVaR_\alpha^{F_Y^L} \leq CVaR_\alpha^{F_X} \leq CVaR_\alpha^{F_Y^U},
	\end{equation}
	where $CVaR_\alpha^{F_Y^L}$, $CVaR_\alpha^{F_X}$, and $CVaR_\alpha^{F_Y^U}$ denote the CVaR at level $\alpha$ corresponding to the distributions $F_Y^L$, $F_X$, and $F_Y^U$, respectively.
	
	Let $a_{\min} = \min(a_X, a_Y)$ and $b_{\min} = \min(b_X, b_Y)$. Define the interval $[b_{\min}, b_X)$ to be $\emptyset$ if $b_{\min} = b_X$, and equal to $[b_Y, b_X)$ otherwise. Analogously, define $[a_{\min}, a_X)$ in the same manner. We then define upper and lower bounds for $F_Y$ as follows (see Figure \ref{fig:CDF_bounds}): \begin{equation}
		F_Y^U(y)=\begin{cases}
			0 & y< max(a_X,q_{1-\epsilon}^Y) \\
			min(F_Y(y)-\epsilon, 1-\epsilon) & y\in [max(a_X,q_{1-\epsilon}^Y), b_{min}) \\
			1-\epsilon & y\in [b_{min}, b_X) \\
			1 & y\geq max(b_X, b_Y),
		\end{cases}
	\end{equation}
	\begin{equation}
		F_Y^L(y)=\begin{cases}
			0 & y<a_{min}  \\
			\epsilon & y\in [a_{min}, a_Y) \\
			min(F_Y(y)+\epsilon, 1) & y\in [a_Y, min(q_{\epsilon}^Y, b_X)) \\
			1 & y\geq min(q_\epsilon^Y, b_X).
		\end{cases}
	\end{equation}
	As a first step, we verify that $F_Y^L$ and $F_Y^U$ are valid CDFs and that they satisfy $F_Y^L \leq F_X \leq F_Y^U$. To establish that $F$ is a CDF, it suffices to verify the following properties:
	\begin{enumerate}
		\item $F$ is non-decreasing;
		\item $F : \mathbb{R} \to [0,1]$ with $\lim_{x \to \infty} F(x) = 1$ and $\lim_{x \to -\infty} F(x) = 0$;
		\item $F$ is right-continuous.
	\end{enumerate}
	\textbf{Proof that $F_Y^L$ is a CDF:} \begin{enumerate}
		\item On the interval $[a_Y, min(q_{1-\epsilon}^Y, b_X))$, the function $F_Y^L$ is monotone increasing, since $F_Y$ is monotone increasing by virtue of being a CDF. Outside this interval, $F_Y^L$ is constant and consequently preserves monotonicity.
		
		\item By definition, $lim_{y\rightarrow \infty}F_Y^L(y)=1$ and $lim_{y\rightarrow -\infty}F_Y^L(y)=0$. We need to show that $F_Y^L$ is bounded between 0 and 1. By its definition, $F_Y^L$ is bounded between 0 and 1.
		
		\item Within the interval $[a_Y, min(q_{\epsilon}^Y, b_X))$, $F_Y^L$ is right-continuous since $F_Y$ is right-continuous as a CDF. Outside this interval, $F_Y^L$ is constant and hence also right-continuous.
	\end{enumerate}
	\textbf{Proof that $F_Y^U$ is a CDF:} \begin{enumerate}
		\item On the interval $[max(a_X,q_{1-\epsilon}^Y), b_{min})$, the function $F_Y^U$ is monotone increasing, since $F_Y$ is monotonically increasing by virtue of being a CDF. Outside this interval, $F_Y^U$ is constant and consequently preserves monotonicity.
		\item By definition, $lim_{y\rightarrow \infty}F_Y^U(y)=1$ and $lim_{y\rightarrow -\infty}F_Y^U(y)=0$. By its definition, $F_Y^U$ is bounded between 0 and 1.
		\item Within the interval $[max(a_X,q_{1-\epsilon}^Y), b_{min})$, $F_Y^U$ is right-continuous since $F_Y$ is right-continuous as a CDF. Outside this interval, $F_Y^U$ is constant and hence also right-continuous.
	\end{enumerate}
	\textbf{Proof that $F_Y^L \leq F_X \leq F_Y^U$:} To establish that $F_Y^L \leq F_X \leq F_Y^U$, we need to show that for all $z \in \mathbb{R}$, $F_Y^L(z) \geq F_X(z) \geq F_Y^U(z)$. Let $z\in \mathbb{R}$. \begin{itemize}
		\item If $z<a_{min}$ then $F_Y^L(z)=0 = F_X(z)$.
		\item If $z\in [a_{min}, a_Y)$ then $F_X(y)=F_X(y)-F_Y(y)+F_Y(y)\leq |F_X(y)-F_Y(y)| + F_Y(y)\leq \epsilon + F_Y(y)=\epsilon=F_Y^L(y)$.
		\item If $z\in [a_Y, min(q_\epsilon^Y, b_X))$ we assume that $F_Y(z)+\epsilon\leq 1$ because otherwise $F_Y^L(z)=1$ and the inequality holds. $F_Y^L(z)=F_Y(z) + \epsilon = F_Y(z)-F_X(z)+F_X(z)+\epsilon \geq F_X(z) - |F_Y(z)-F_X(z)| + \epsilon\geq F_X(z)$.
		\item If $z\geq min(q_\epsilon^Y, b_X)$ then $F_Y^L(z)=1\geq F_X(z)$.
	\end{itemize} 
	and therefore $F_Y^L\leq F_X$. Note that the last proof holds when $b_X=\infty$, making $F_Y^L$ a valid CDF when the support of $X$ is not bounded from above.
	\begin{itemize}
		\item If $z<max(a_X,q_{1-\epsilon}^Y)$ then $F_Y^U(z)=0\leq F_X(z)$.
		\item If $z\in [max(a_X,q_{1-\epsilon}^Y), b_{min})$ then $F_Y^U(z)\leq F_Y(z)-\epsilon=F_Y(z)-F_X(z)+F_X(z)-\epsilon \leq |F_Y(z)-F_X(z)|+F_X(z)-\epsilon \leq \epsilon+F_X(z)-\epsilon= F_X(z)$.
		\item If $z \in [b_{min}, b_X)$ then $F_Y^U(z)=1-\epsilon=F_Y(z)-\epsilon=F_Y(z)-\epsilon+F_X(z)-F_X(z)\leq F_X(z)-\epsilon+|F_Y(z)-F_X(z)|\leq F_X(z)$
	\end{itemize}
	and therefore $F_X\leq F_Y^U$. Note that the last proof holds when $a_X=-\infty$, making $F_Y^U$ a valid CDF when the support of $X$ is not bounded from below.
	
	Next, we derive bounds for the CVaR associated with $F_Y^U$ and $F_Y^L$.
	
	\textbf{Upper bound for $CVaR_\alpha^{F_Y^U}$:} From \cite{acerbi2002coherence}, CVaR is equal to an integral of the VaR
	\begin{equation}
		\begin{aligned}
			&CVaR_\alpha^{F_Y^U}=\frac{1}{\alpha}\int_{1-\alpha}^1 \inf \{y\in \mathbb{R}:F_Y^U(y)\geq \tau\} d\tau \\
			&=\frac{1}{\alpha}\int_{1-\alpha+\epsilon}^{1+\epsilon} \inf \{y\in \mathbb{R}:F_Y^U(y)\geq \tau -\epsilon\} d\tau
		\end{aligned}
	\end{equation}  
	If $\alpha \leq \epsilon$, then $1 - \alpha + \epsilon \geq 1$, rendering the bound in the preceding equation trivial, as it attains the maximum value of the support of both $X$ and $Y$.
	\begin{equation}\label{eq:helper1}
		\begin{aligned}
			&\frac{1}{\alpha}\int_{1-\alpha+\epsilon}^{1+\epsilon} \inf \{y\in \mathbb{R}:F_Y^U(y)\geq \tau -\epsilon\} \\
			&\leq \frac{1}{\alpha}\int_{1-\alpha+\epsilon}^{1+\epsilon} max(b_X, b_Y)d\tau = max(b_X, b_Y).
		\end{aligned}
	\end{equation}
	That is, $CVaR_\alpha(X)\leq max(b_X, b_Y)$.
	
	If $\alpha > \epsilon$, then $1 - \alpha + \epsilon < 0$, and the integral may be decomposed into one term that is trivially bounded by the maximum of the support and another term that can be computed explicitly
	\begin{equation}
		\begin{aligned}
			&CVaR_\alpha^{F_Y^U}=\frac{1}{\alpha}\int_{1-\alpha+\epsilon}^{1+\epsilon} \inf \{y\in \mathbb{R}:F_Y^U(y)\geq \tau -\epsilon\} d\tau \\
			&=\frac{1}{\alpha}[\underbrace{\int_{1}^{1+\epsilon} \inf \{y\in \mathbb{R}:F_Y^U(y)\geq \tau -\epsilon\} d\tau}_{\triangleq A_1} \\
			&+\underbrace{\int_{1-\alpha+\epsilon}^{1} \inf \{y\in \mathbb{R}:F_Y^U(y)\geq \tau -\epsilon\} d\tau}_{\triangleq A_2}]
		\end{aligned}
	\end{equation}
	The term $A_1$ is bounded, in a manner analogous to \eqref{eq:helper1}, by $\epsilon \max(b_X, b_Y)$, which is essentially the tightest bound attainable given the definition of $F_Y^U$. The term $A_2$ can be expressed in terms of the CVaR of $Y$, evaluated at a shifted confidence level. Specifically, the variable $\tau - \epsilon$ in $A_2$ lies within the interval $[1 - \alpha, 1 - \epsilon]$. Over this range, for all $y \in [\max(a_X, q_{1-\epsilon}^Y), b_{\min}]$, the inequality $F_Y^U(y) \leq F_Y(y) - \epsilon$ holds. This follows because if $q_{1-\epsilon}^Y \geq a_X$, then by definition $F_Y^U(y) = min(F_Y(y) - \epsilon, 1-\epsilon)$, and if $a_X > q_{1-\epsilon}^Y$, then $F_Y^U(y)=0 \leq F_Y(y) - \epsilon$ for $y\in [q_{1-\epsilon}^Y, a_X]$ and again $F_Y^U(y) = min(F_Y(y) - \epsilon, 1-\epsilon)$ for $y\in [a_X, b_{min}]$.
	\begin{equation}
		\begin{aligned}
			&A_2 \leq \int_{1-\alpha+\epsilon}^1 \inf \{y\in \mathbb{R}:F_Y(y) - \epsilon \geq \tau -\epsilon\}d\tau \\
			&=\int_{1-(\alpha-\epsilon)}^1 \inf \{y\in \mathbb{R}:F_Y(y) \geq \tau\}d\tau \\
			&=(\alpha-\epsilon)CVaR_{\alpha-\epsilon}(Y).
		\end{aligned}
	\end{equation}
	Note that the confidence level is valid, as $\alpha, \epsilon \in (0,1)$ and $\alpha > \epsilon$ imply $\alpha-\epsilon \in (0, 1)$. By combining the previous two expressions, we obtain a single bound:
	\begin{equation}
		CVaR_\alpha^{F_Y^U} \leq \epsilon\frac{max(b_X,b_Y)}{\alpha} + (1 - \frac{\epsilon}{\alpha})CVaR_{\epsilon-\alpha}(Y).
	\end{equation}
	In the case where $a_X = -\infty$, we have $\max(a_X, q_{1-\epsilon}^Y) = q_{1-\epsilon}^Y$, and the definition of $F_Y^U$ no longer involves $a_X$. Therefore, the inequality remains valid in the case of $a_X = -\infty$.
	
	\textbf{Lower bound for $CVaR_\alpha^{F_Y^L}$:} From \cite{acerbi2002coherence}, CVaR is equal to an integral of the VaR
	\begin{equation}
		\begin{aligned}
			&CVaR_\alpha^{F_Y^L}=\frac{1}{\alpha}\int_{1-\alpha}^1 \inf \{y\in \mathbb{R}:F_Y^L(y)\geq \tau\} d\tau \\
			&=\frac{1}{\alpha}\int_{1-(\epsilon+\alpha)}^{1-\epsilon} \inf \{y\in \mathbb{R}:F_Y^L(y)\geq \tau +\epsilon\} d\tau.
		\end{aligned}
	\end{equation}
	It holds that $F_Y^L(y)\leq F_Y(y)+\epsilon$ and therefore
	\begin{equation}
		\begin{aligned}
			&CVaR_\alpha^{F_Y^L} \geq \frac{1}{\alpha}\int_{1-(\epsilon+\alpha)}^{1-\epsilon} \inf \{y\in \mathbb{R}:F_Y(y)+\epsilon \geq \tau +\epsilon\} d\tau \\
			&=\frac{1}{\alpha}\int_{1-(\epsilon+\alpha)}^{1-\epsilon} \inf \{y\in \mathbb{R}:F_Y(y)\geq \tau\} d\tau \\
			&=\frac{1}{\alpha}[\int_{1-(\epsilon+\alpha)}^{1} \inf \{y\in \mathbb{R}:F_Y(y)\geq \tau\} d\tau \\
			&-\int_{1-\epsilon}^{1} \inf \{y\in \mathbb{R}:F_Y(y)\geq \tau\} d\tau] \\
			&=\frac{1}{\alpha}[(\alpha+\epsilon)CVaR_{\alpha+\epsilon}(Y)-\epsilon CVaR_{\epsilon}(Y)].
		\end{aligned}
	\end{equation}
	In the case where $\epsilon+\alpha > 1$,
	\begin{equation}
		\begin{aligned}
			&CVaR_\alpha^{F_Y^L}=\frac{1}{\alpha}\int_{1-(\epsilon+\alpha)}^{1-\epsilon} \inf \{y\in \mathbb{R}:F_Y^L(y)\geq \tau +\epsilon\} d\tau \\
			&=\frac{1}{\alpha}[\underbrace{\int_{0}^{1-\epsilon} \inf \{y\in \mathbb{R}:F_Y^L(y)\geq \tau +\epsilon\} d\tau}_{\triangleq A_3} \\
			&+\underbrace{\int_{1-(\epsilon+\alpha)}^{0} \inf \{y\in \mathbb{R}:F_Y^L(y)\geq \tau +\epsilon\} d\tau}_{\triangleq A_4}].
		\end{aligned}
	\end{equation}
	
	\begin{equation}
		\begin{aligned}
			&A_3 \geq \int_{0}^{1-\epsilon} \inf \{y\in \mathbb{R}:F_Y(y)\geq \tau\} d\tau \\
			&=\int_{0}^{1} \inf \{y\in \mathbb{R}:F_Y(y)\geq \tau\} d\tau \\
			&-\int_{1-\epsilon}^{1} \inf \{y\in \mathbb{R}:F_Y(y)\geq \tau\} d\tau \\
			&=\mathbb{E}[Y] - \epsilon CVaR_\epsilon(Y)
		\end{aligned}
	\end{equation}
	In the case of $A_4$, $\tau$ ranges from $1-\alpha$ to $\epsilon$, and therefore 
	\begin{equation}
		\begin{aligned}
			&A_4=\int_{1-\alpha}^{\epsilon} \inf \{y\in \mathbb{R}:F_Y^L(y)\geq \tau\} d\tau \geq (\alpha+\epsilon-1)a_{min}
		\end{aligned}
	\end{equation}
	By combining the last equations to one bound we get \begin{equation}
		\begin{aligned}
			CVaR_\alpha(X)&\geq CVaR_\alpha^{F_Y^L} \\
			&\geq \mathbb{E}[Y] - \epsilon CVaR_\epsilon(Y) + (\alpha+\epsilon-1)a_{min}.
		\end{aligned}
	\end{equation}
\end{proof}

\begin{theorem}\label{thm:cvar_bound_convergence}
	Given the definition of $X,Y,\epsilon$ and $\alpha$ as in Theorem \ref{thm:cvar_bound_v2}, the lower and upper bounds of Theorem \ref{thm:cvar_bound_v2} converge to $CVaR_\alpha(X)$ as $\epsilon \rightarrow 0$.
\end{theorem}
\begin{proof}
	Let $f$ and $g$ be continuous functions such that $f(0)$ and $g(0)$ are finite. Then
	$\lim_{x \to 0} f(x) g(x) = \big( \lim_{x \to 0} f(x) \big) \big( \lim_{x \to 0} g(x) \big)$.
	This property will be used throughout the proof.
	
	Assume that $F_Y(z) - F_X(z) \leq \epsilon$ for all $z \in \mathbb{R}$. Since we consider the limit as $\epsilon \to 0$, we further assume that $\alpha > \epsilon$ when computing the bound. Noting that CVaR is continuous with respect to the confidence level, it follows that $\lim_{\epsilon \to 0} \mathrm{CVaR}_{\alpha - \epsilon}(Y) = \mathrm{CVaR}_{\alpha}(Y)$.
	\begin{equation}
		\begin{aligned}
			&\lim_{\epsilon\rightarrow 0} \frac{\epsilon}{\alpha}max(b_X,b_Y) + (1 - \frac{\epsilon}{\alpha})CVaR_{\alpha-\epsilon}(Y) \\
			&=\lim_{\epsilon \rightarrow 0} (1 - \frac{\epsilon}{\alpha})CVaR_{\alpha-\epsilon}(Y) \\
			&=\lim_{\epsilon \rightarrow 0} 1 - \frac{\epsilon}{\alpha} \lim_{\epsilon\rightarrow 0} CVaR_{\alpha-\epsilon}(Y)=CVaR_\alpha(Y)
		\end{aligned}
	\end{equation}
	
	If $F_X(z) - F_Y(z) \leq \epsilon$ for all $z \in \mathbb{R}$, then, since we consider the limit as $\epsilon \to 0$, we assume in the derivation of the bound that $\alpha + \epsilon \leq 1$.
	\begin{equation}
		\begin{aligned}
			&\lim_{\epsilon\rightarrow 0} (1+\frac{\epsilon}{\alpha})CVaR_{\alpha+\epsilon}(Y)-\frac{\epsilon}{\alpha} CVaR_{\epsilon}(Y) \\
			&=\lim_{\epsilon\rightarrow 0} 1+\frac{\epsilon}{\alpha} \lim_{\epsilon\rightarrow 0} CVaR_{\alpha+\epsilon}(Y) \\
			&-\lim_{\epsilon\rightarrow 0} \frac{\epsilon}{\alpha} \lim_{\epsilon \rightarrow 0} CVaR_{\alpha+\epsilon}(Y) \\
			&=CVaR_\alpha(Y).
		\end{aligned}
	\end{equation}
\end{proof}

\begin{theorem}
	\label{thm:ecdf_cvar_bound}
	Let $X$ a random variable, $\alpha\in (0, 1],\delta \in (0, 1), \epsilon=\sqrt{\ln(1/\delta)/(2n)}$. Let $X_1, \dots, X_n \overset{iid}{\sim} F_X$ be random variables that define the ECDF $\hat{F}_X$. 
	\begin{enumerate}
		\item \textbf{Upper Bound:} If $P(X\leq b)=1$, then 
		\begin{enumerate}
			\item If $\alpha > \epsilon$ then $P(CVaR_\alpha(X) \leq (1-\frac{\epsilon}{\alpha})C_{\alpha-\epsilon}^{\hat{F}_X} + \frac{\epsilon}{\alpha}b)> 1-\delta$.
			
			\item If $\alpha\leq \epsilon$ then $CVaR_\alpha(X)\leq b$
		\end{enumerate}
		
		\item \textbf{Lower Bound:} If $P(X\geq a)=1$, then 
		\begin{enumerate}
			\item If $\alpha+\epsilon < 1$, then $P(CVaR_\alpha(X)\geq (1+\frac{\epsilon}{\alpha})C_{\alpha+\epsilon}^{\hat{F}_X}-\frac{\epsilon}{\alpha}C_{\epsilon}^{\hat{F}_X})> 1-\delta$.
			
			\item If $\alpha+\epsilon \geq 1$, then $P(CVaR_\alpha(X) \geq \frac{1}{\alpha}[(\alpha+\epsilon - 1)a + E_{\hat{F}_X}[X] - \epsilon C_\epsilon^{\hat{F}_X}]) > 1-\delta$.
		\end{enumerate}
	\end{enumerate}
\end{theorem}
\begin{proof}
	From DKW inequality the following inequalities can be derived \cite{pmlr-v97-thomas19a}
	\begin{equation}
		\Pr \left( \sup_{x \in \mathbb{R}} \big( \hat{F}(x) - F(x) \big) 
		\leq \sqrt{\frac{\ln(1/\delta)}{2n}} \right) \geq 1 - \delta,
	\end{equation}
	\begin{equation}
		\Pr \left( \sup_{x \in \mathbb{R}} \big( \hat{F}(x) - F(x) \big) 
		\geq \sqrt{\frac{\ln(1/\delta)}{2n}} \right) \geq 1 - \delta.
	\end{equation}
	Let $\epsilon=\sqrt{\frac{\ln(1/\delta)}{2n}}$. By using Theorem \ref{thm:cvar_bound_v2} we get the following. 
	
	If $\alpha< \epsilon$ then \begin{equation}
		\begin{aligned}
			&P(CVaR_\alpha(X) \leq \frac{\alpha-\epsilon}{\alpha}C_{\alpha-\epsilon}^{\hat{F}_X} + \frac{\epsilon}{\alpha}b) \\ 
			&= \underbrace{P(C_\alpha(X) \leq \frac{\alpha-\epsilon}{\alpha}C_{\alpha-\epsilon}^{\hat{F}_X} + \frac{\epsilon}{\alpha}b\Bigr{|} \sup_{z\in\mathbb{R}}(\hat{F}_X-F_X)\leq \epsilon)}_{=1} \\ 
			&\times \underbrace{P(\sup_{z\in\mathbb{R}}(\hat{F}_X-F_X)\leq \epsilon)}_{>1-\delta} \\
			&+ \underbrace{P(C_\alpha(X) \leq \frac{\alpha-\epsilon}{\alpha}C_{\alpha-\epsilon}^{\hat{F}_X} + \frac{\epsilon}{\alpha}b\Bigr{|} \sup_{z\in\mathbb{R}}(\hat{F}_X-F_X)> \epsilon)}_{\geq 0} \\
			&\times \underbrace{P(\sup_{z\in\mathbb{R}}(\hat{F}_X-F_X)> \epsilon)}_{\geq 0} \\
			&> 1-\delta.
		\end{aligned}
	\end{equation}
	Observe that, conditional on the event 
	\[
	\sup_{z \in \mathbb{R}} \big( \hat{F}_X(z) - F_X(z) \big) \leq \epsilon,
	\]
	the bound in Theorem~\ref{thm:cvar_bound_v2} holds deterministically. Consequently, the probability that
	\[
	C_\alpha(X) \leq \frac{\alpha - \epsilon}{\alpha} C_{\alpha - \epsilon}^{\hat{F}_X} + \frac{\epsilon}{\alpha} b
	\]
	holds, given 
	\[
	\sup_{z \in \mathbb{R}} \big( \hat{F}_X(z) - F_X(z) \big) \leq \epsilon,
	\]
	is equal to one. The same observation is necessary through the rest of the proof in a similar manner.
	If $\alpha+\epsilon < 1$, then \begin{equation}
		\begin{aligned}
			&P(C_\alpha(X)\geq (1+\frac{\epsilon}{\alpha})C_{\alpha+\epsilon}^{\hat{F}_X}-\frac{\epsilon}{\alpha}C_{\epsilon}^{\hat{F}_X}) \\
			&= \underbrace{P(C_\alpha(X)\geq (1+\frac{\epsilon}{\alpha})C_{\alpha+\epsilon}^{\hat{F}_X}-\frac{\epsilon}{\alpha}C_{\epsilon}^{\hat{F}_X}\Bigr{|} \sup_{z\in\mathbb{R}}(F_X-\hat{F}_X)\leq \epsilon)}_{=1} \\ &\underbrace{P(\sup_{z\in\mathbb{R}}(F_X-\hat{F}_X)\leq \epsilon)}_{>1-\delta} \\
			&+\underbrace{P(C_\alpha(X)\geq (1+\frac{\epsilon}{\alpha})C_{\alpha+\epsilon}^{\hat{F}_X}-\frac{\epsilon}{\alpha}C_{\epsilon}^{\hat{F}_X}\Bigr{|} \sup_{z\in\mathbb{R}}(F_X-\hat{F}_X) > \epsilon)}_{\geq 0} \\
			&\times \underbrace{P(\sup_{z\in\mathbb{R}}(F_X-\hat{F}_X) > \epsilon)}_{\geq 0} \\
			&> 1-\delta
		\end{aligned}
	\end{equation}
	If $\alpha+\epsilon \geq 1$, then \begin{equation}
		\begin{aligned}
			&P(C_\alpha(X) \geq \frac{1}{\alpha}[(\alpha+\epsilon - 1)a + E_{\hat{F}_X}[X] - \epsilon C_\epsilon^{\hat{F}_X}]) \\ 
			&=\underbrace{P(C_\alpha(X) \geq \frac{1}{\alpha}[(\alpha+\epsilon - 1)a + E_{\hat{F}_X}[X]}_{=1} \\
			&\underbrace{- \epsilon C_\epsilon^{\hat{F}_X}]\Bigr{|} \sup_{z\in\mathbb{R}}(F_X-\hat{F}_X)\leq \epsilon)}_{=1} \\ 
			&\times \underbrace{P(\sup_{z\in\mathbb{R}}(F_X-\hat{F}_X) \leq \epsilon)}_{>1-\delta} \\
			&+\underbrace{P(C_\alpha(X) \geq \frac{1}{\alpha}[(\alpha+\epsilon - 1)a + E_{\hat{F}_X}[X]}_{\geq 0} \\
			&\underbrace{- \epsilon C_\epsilon^{\hat{F}_X}]\Bigr{|} \sup_{z\in\mathbb{R}}(F_X-\hat{F}_X) > \epsilon)}_{\geq 0} \\ 
			&\times \underbrace{P(\sup_{z\in\mathbb{R}}(F_X-\hat{F}_X) > \epsilon)}_{\geq 0} \\
			&> 1-\delta
		\end{aligned}
	\end{equation}
\end{proof}

\begin{theorem}\label{thm:tight_cvar_lower_bound}
	(Tighter CVaR Lower Bound) Let $\alpha \in (0,1)$, $X$ and $Y$ be random variables. Define the a random variable $Y^L$ such that $F_{Y^L}(y) \triangleq min(1, F_Y(y) + g(y))$ for $g:\mathbb{R}\rightarrow [0, \infty)$. Assume 
	$lim_{x \rightarrow -\infty}g(x)=0$, $g$ is continuous from the right and monotonic increasing.
	If $\forall x\in \mathbb{R}, F_X(x)\leq F_Y(x)+g(x)$, then $F_{Y^L}$ is a CDF and $CVaR_\alpha(Y^L)\leq CVaR_\alpha(X).$
\end{theorem}
\begin{proof}
	In order to prove that $F_{Y^L}$ is a CDF we need to prove that $F_{Y^L}$ is:
	\begin{enumerate}
		\item Monotonic increasing
		\item $F:\mathbb{R}:\rightarrow [0,1]$, $\lim_{x\rightarrow \infty}F_{Y^L}(x)=1$, $\lim_{x\rightarrow -\infty}F_{Y^L}(x)=0$
		\item Continuous from the right
	\end{enumerate}
	
	\textbf{Monotonic increasing:}
	Note that for every $f_i:\mathbb{R}\rightarrow \mathbb{R}, i=1,2$ that are monotonic increasing, $f_1(f_2(x)))$ is also monotonic increasing in x. Denote $f(x):=min(x,1)$ and $f_2(x):=F_Y(x)+g(x)$. $F_Y(x)$ is a CDF and therefore monotonic increasing, so $f_2(x)$ is monotonic increasing as a sum of monotonic increasing functions. $f_1$ is also monotonic increasing, and $F_{Y^L}(x)=f_1(f_2(x))$. Therefore $F_{Y^L}(x)$ is monotonic increasing. \\
	
	\textbf{Limits:} \begin{equation}
		\begin{aligned}
			1&\geq\lim_{x\rightarrow \infty} F_{Y^L}(x)=\lim_{x\rightarrow \infty} min(1,F_Y(x)+g(x)) \\
			&\geq \lim_{x\rightarrow \infty}min(1,F_Y(x))=\lim_{x\rightarrow \infty} F_Y(x)=1
		\end{aligned}
	\end{equation}
	and therefore $\lim_{x\rightarrow \infty} F_{Y^L}(x)=1$.
	$$\begin{aligned}
		0&\leq \lim_{x\rightarrow -\infty} F_{Y^L}(x)=\lim_{x\rightarrow -\infty} min(1,F_Y(x)+g(x)) \\
		&\leq \lim_{x\rightarrow -\infty} F_Y(x)+g(x) 
		=\lim_{x\rightarrow -\infty} F_Y(x)+\lim_{x\rightarrow -\infty} g(x) \\
		&=0	
	\end{aligned}$$
	and therefore $\lim_{x\rightarrow -\infty} F_{Y^L}(x)=0$. By definition $\forall x\in \mathbb{R},F_{Y^L}(x)\leq 1$, and $\forall x\in \mathbb{R},F_{Y^L}(x)\geq 0$ because both g and $F_{Y^L}$ are non negative functions. \\
	
	\textbf{Continuity from the right:}
	$F_Y$ is continuous from the right because it is a CDF, and therefore $F_{Y^L}$ is continuous from the right as a sum of continuous from the right functions. \\
	Thus, $F_{Y^L}$ is a CDF. \\
	
	\textbf{Bound proof:}
	If $Y^L\leq X$, then $CVaR_\alpha(Y^L)\leq CVaR_\alpha(X)$ because CVaR is a coherent risk measure. Note that 
	if $F_Y(x) + g(x)<1$ then
	$$F_X(x)\leq F_Y(x) + g(x)=F_{Y^L}(x),$$
	and if $F_Y(x) + g(x)\geq 1$, $1=F_{Y^L}(x)\geq F_X(x)$. Therefore $Y^L\leq X$.
\end{proof}

\begin{theorem}\label{thm:lower_cvar_bound_using_density_bound}
	Let $\alpha \in (0,1)$, $X$ and $Y$ random variables. Define $h:\mathbb{R}\rightarrow [0,\infty)$ to be a continuous function, $g(z):=\int_{-\infty}^z h(x)dx$ and $Y^L$ to be a random variable such that $F_{Y^L}(y):=\min(1, F_{Y}(y) + g(y))$. If $\lim_{z\rightarrow -\infty} g(z)=0$ and $\forall x\in \mathbb{R},f_x(x)\leq f_y(z) + h(x)$, then $F_{Y^L}$ is a CDF and $CVaR_\alpha(Y^L)\leq CVaR_\alpha(X)$.
\end{theorem}
\begin{proof}
	We will show the g satisfies the properties of Theorem \ref{thm:tight_cvar_lower_bound}, and therefore this theorem holds. We need to prove that \begin{enumerate}
		\item $\lim_{z\rightarrow -\infty} g(z)=0$
		\item g is continuous from the right.
		\item g is monotonic increasing.
		\item $F_X(y)\leq F_Y(y) + g(y)$
	\end{enumerate}
	It is given in the theorem's assumptions that $\lim_{z\rightarrow -\infty} g(z)=0$, so (1) holds. h is non negative and therefore $g$ is monotonic increasing, so (3) holds. $g$ is continuous if its derivative exists for all $z\in \mathbb{R}$. Let $z\in \mathbb{R}$ and $a<z$.
	\begin{equation}
		\begin{aligned}
			\frac{d}{dz}g(z) &= \frac{d}{dz}\int_{-\infty}^z h(x)dx = \frac{d}{dz}[\int_{-\infty}^a h(x)dx + \int_{a}^z h(x)dx] \\
			&=\frac{d}{dz}\int_{a}^z h(x)dx = h(z)
		\end{aligned}
	\end{equation}
	where the third equality holds because $\int_{-\infty}^a h(x)dx = g(a)$ is a constant that does not depend on z. The last equality holds from the fundamental theorem of calculus because $h$ is continuous. Finally, (4) holds because \begin{equation}
		\begin{aligned}
			F_X(y) &\triangleq \int_{-\infty}^y f_x(x)dx\leq \int_{-\infty}^y f_y(x) + h(x)dx \\
			&\triangleq F_Y(y) + g(y).
		\end{aligned}
	\end{equation}
\end{proof}

\begin{corollary}\label{thm:x_concentration_bound_asymptotic_convergence}
	Let $X$ be a random variable, $\alpha\in (0, 1],\delta \in (0, 1), \epsilon=\sqrt{\ln(1/\delta)/(2n)}, a\in \mathbb{R}, b\in \mathbb{R}, \eta>0$. Let $X_1, \dots, X_n \overset{iid}{\sim} F_X$ be random variables that define the ECDF $\hat{F}_X$. Denote by $U(n)$ and $L(n)$ the upper and lower bounds respectively from Theorem \ref{thm:ecdf_cvar_bound}, where $n$ is the number of samples, then 
	\begin{enumerate}
		\item If $P(X\leq b)=1$, then $\lim_{n\rightarrow \infty} U(n)\overset{a.s}{=}CVaR_\alpha(X)$.
		\item If $P(X\geq a)=1$, then $\lim_{n\rightarrow \infty} L(n)\overset{a.s}{=}CVaR_\alpha(X)$.
	\end{enumerate}
	where a.s denotes almost sure convergence.
\end{corollary}
\begin{proof}
	\begin{equation}
		\begin{aligned}
			&\lim_{n \to \infty} U(n)\overset{a.s}{=} \lim_{n \to \infty} (1-\frac{\epsilon}{\alpha})C_{\alpha-\epsilon}^{\hat{F}_X} + \frac{\epsilon}{\alpha}b \\
			&\overset{a.s}{=}\lim_{n \to \infty} C_{\alpha-\epsilon}^{\hat{F}_X}=\lim_{n \to \infty} C_{\alpha}^{\hat{F}_X} = C_{\alpha}^{F_X}
		\end{aligned}
	\end{equation}
	where the first and second equalities follow from the fact that $\epsilon \to 0$ as $n \to \infty$; the third equality holds by the continuity of CVaR with respect to $\alpha$; and the fourth equality holds by the almost sure convergence of the empirical CVaR of i.i.d. samples to the true CVaR. For the same reason,
	\begin{equation}
		\begin{aligned}
			&\lim_{n\rightarrow \infty} L(n)\overset{a.s}{=} \lim_{n\rightarrow \infty} (1+\frac{\epsilon}{\alpha})C_{\alpha+\epsilon}^{\hat{F}_X}-\frac{\epsilon}{\alpha}C_{\epsilon}^{\hat{F}_X} \\
			&\overset{a.s}{=} \lim_{n\rightarrow \infty} C_{\alpha+\epsilon}^{\hat{F}_X}
			\overset{a.s}{=} \lim_{n\rightarrow \infty} C_{\alpha}^{\hat{F}_X}
			\overset{a.s}{=} \lim_{n\rightarrow \infty} C_{\alpha}^{F_X}.
		\end{aligned}
	\end{equation}
\end{proof}


\section{Our concentration bounds vs \cite{pmlr-v97-thomas19a}}
\begin{figure}[htbp]
	\centering
	\includegraphics[width=\textwidth]{cvar_bounds_comparison.png}
	\caption{
		Comparison of concentration inequalities for CVaR.
		% Replace this with your actual caption.
	}
	\label{fig:cvar_bounds_comparison}
\end{figure}


\end{document}
